%% BLOCK 2 -- Machine Learning Crash Course.
%% Fragment of ai_foundations.tex; not a standalone document.
%% Source: developers.google.com/machine-learning/crash-course, all twelve
%% modules and their unit pages. The module order is the course's own
%% recommended order: ML models, then data, then advanced models, then
%% real-world ML.

\actpage{2}{Machine Learning Crash Course}{Twelve self-contained modules: how to build regression and classification models, how to work with data, and what changes in the real world.}

\begin{frame}{Twelve modules, four groups}
  \begin{withmargin}
    \begin{tdmain}
      \footnotesize
      \begin{tabular}{@{}p{0.26\linewidth}p{0.68\linewidth}@{}}
        \toprule
        \textbf{Group} & \textbf{Modules}\\
        \midrule
        ML models & Linear regression, logistic regression, classification\\[0.3em]
        Data & Numerical data, categorical data, datasets and overfitting\\[0.3em]
        Advanced models & Neural networks, embeddings, large language models\\[0.3em]
        Real-world ML & Production ML systems, AutoML, ML fairness\\
        \bottomrule
      \end{tabular}
    \end{tdmain}
    \begin{tdmargin}
      {\color{tdfg}\textbf{Self-contained.}}
      \par\medskip
      Each module is self-contained, so with prior experience you can skip
      directly to the topics you want. If you are new to ML, complete the
      modules in the order shown.
      \par\medskip
      \gterm{Data} gets three of the twelve modules. ML practitioners spend far
      more time evaluating, cleaning, and transforming data than building
      models.
    \end{tdmargin}
  \end{withmargin}
  \slidesources{GoogleMLCC2025}
\end{frame}

% ======================================================================
\sectionsubtitle{Finding the relationship between features and a label.}
\section[Linear regression]{Linear Regression}

\begin{frame}{\modulehead{Linear regression}{80 minutes}}
  \begin{withmargin}
    \begin{tdmain}
      Linear regression is a statistical technique used to find the relationship
      between variables. In an ML context, linear regression finds the
      relationship between \gterm{features} and a \gterm{label}.
      \vskip0.7em
      \footnotesize
      Suppose we want to predict a car's fuel efficiency in miles per gallon
      based on how heavy the car is:
      \vskip0.4em
      \begin{tabular}{@{}rrrrrrr@{}}
        \toprule
        \multicolumn{7}{@{}l}{Pounds in 1000s (feature)}\\
        3.5 & 3.69 & 3.44 & 3.43 & 4.34 & 4.42 & 2.37\\
        \midrule
        \multicolumn{7}{@{}l}{Miles per gallon (label)}\\
        18 & 15 & 18 & 16 & 15 & 14 & 24\\
        \bottomrule
      \end{tabular}
    \end{tdmain}
    \begin{tdmargin}
      {\color{tdfg}\textbf{Objectives.}}
      \par\medskip
      Explain a loss function and how it works. Define and describe how gradient
      descent finds the optimal model parameters. Describe how to tune
      hyperparameters to efficiently train a linear model.
      \par\medskip
      As a car gets heavier, its miles per gallon rating generally decreases.
    \end{tdmargin}
  \end{withmargin}
\end{frame}

\begin{frame}{The linear regression equation}
  \begin{withmargin}
    \begin{tdmain}
      \[ y' = b + w_1 \gterm{x_1} \]
      \vskip0.4em
      \footnotesize
      In algebraic terms the model is $y = mx + b$. In ML we rename the parts:
      \vskip0.3em
      \begin{itemize}
        \item $y'$ is the predicted label, the output.
        \item $b$ is the \textbf{bias}, the same concept as the y-intercept.
          Sometimes written $w_0$. It is a parameter of the model and is
          calculated during training.
        \item $w_1$ is the \textbf{weight} of the feature, the same concept as
          the slope $m$. Also a parameter, also calculated during training.
        \item $\gterm{x_1}$ is a feature, the input.
      \end{itemize}
    \end{tdmain}
    \begin{tdmargin}
      {\color{tdfg}\textbf{The fitted car model.}}
      \par\medskip
      Bias is 34, where the line intersects the y-axis. Weight is $-4.6$, the
      slope of the line. So
      \par\smallskip
      $y' = 34 + (-4.6)(x_1)$
      \par\smallskip
      A 4{,}000-pound car has a predicted fuel efficiency of 15.6 miles per
      gallon.
    \end{tdmargin}
  \end{withmargin}
\end{frame}

\begin{frame}{Models with multiple features}
  \begin{withmargin}
    \begin{tdmain}
      A more sophisticated model relies on multiple features, each having a
      separate weight:
      \[ y' = b + w_1x_1 + w_2x_2 + w_3x_3 + w_4x_4 + w_5x_5 \]
      \vskip0.4em
      A model that predicts gas mileage could additionally use engine
      displacement, acceleration, number of cylinders, and horsepower.
      \vskip0.6em
      \footnotesize
      During training, the model calculates the weight and bias that produce the
      best model. The \gterm{bias and weights} are what training updates.
      Predictions are not updated during training, and feature values are part
      of the dataset.
    \end{tdmain}
    \begin{tdmargin}
      As a car's engine gets bigger, its miles per gallon rating generally
      decreases.
      \par\medskip
      As a car's acceleration takes longer, the miles per gallon rating
      generally increases.
      \par\medskip
      Both additional features have a linear relationship to the label.
    \end{tdmargin}
  \end{withmargin}
  \keyterms{bias; feature; label; linear regression; parameter; weight}
\end{frame}

\begin{frame}{Loss}
  \begin{withmargin}
    \begin{tdmain}
      Loss is a numerical metric that describes how wrong a model's predictions
      are. Loss measures the distance between the model's predictions and the
      actual labels. The goal of training a model is to \alertbf{minimize the
      loss}, reducing it to its lowest possible value.
      \vskip0.9em
      Loss focuses on the distance between the values, not the direction. If a
      model predicts 2 but the actual value is 5, we don't care that the loss is
      negative. We care that the distance between the values is 3. Thus all
      methods for calculating loss remove the sign.
    \end{tdmain}
    \begin{tdmargin}
      {\color{tdfg}\textbf{Two ways to remove the sign.}}
      \par\medskip
      Take the absolute value of the difference between the actual value and the
      prediction.
      \par\smallskip
      Square the difference between the actual value and the prediction.
    \end{tdmargin}
  \end{withmargin}
\end{frame}

\begin{frame}{The five types of loss}
  \begin{withmargin}
    \begin{tdmain}
      \footnotesize
      \begin{tabular}{@{}p{0.24\linewidth}p{0.36\linewidth}p{0.32\linewidth}@{}}
        \toprule
        \textbf{Loss type} & \textbf{Definition} & \textbf{Equation}\\
        \midrule
        L1 loss & Sum of the absolute values of the difference
          & $\sum |\text{actual} - \text{predicted}|$\\[0.4em]
        Mean absolute error (MAE) & Average of L1 losses across $N$ examples
          & $\frac{1}{N}\sum |\text{actual} - \text{predicted}|$\\[0.4em]
        L2 loss & Sum of the squared differences
          & $\sum (\text{actual} - \text{predicted})^2$\\[0.4em]
        Mean squared error (MSE) & Average of L2 losses across $N$ examples
          & $\frac{1}{N}\sum (\text{actual} - \text{predicted})^2$\\[0.4em]
        Root mean squared error (RMSE) & Square root of the MSE
          & $\sqrt{\frac{1}{N}\sum (\text{actual} - \text{predicted})^2}$\\
        \bottomrule
      \end{tabular}
    \end{tdmain}
    \begin{tdmargin}
      The functional difference between L1 and L2 is \gterm{squaring}. When the
      difference is large, squaring makes the loss even larger. When the
      difference is small, less than 1, squaring makes it even smaller.
      \par\medskip
      MAE represents the average prediction error. RMSE represents the spread of
      the errors, and is more skewed by larger errors.
    \end{tdmargin}
  \end{withmargin}
\end{frame}

\begin{frame}{Calculating loss}
  \begin{withmargin}
    \begin{tdmain}
      For the car model $y' = 34 + (-4.6)(x_1)$, with weight $-4.6$ and bias
      $34$: a 2{,}370-pound car is predicted to get 23.1 miles per gallon, but
      actually gets 24.
      \vskip0.6em
      \small
      \begin{align*}
        \text{prediction} &= 34 + (-4.6 \cdot 2.37) = 23.1\\
        \text{actual} &= 24\\
        \text{L2 loss} &= (24 - 23.1)^2 = \gterm{0.81}
      \end{align*}
    \end{tdmain}
    \begin{tdmargin}
      The formula uses 2.37 because the graphs are scaled to 1000s of pounds.
      \par\medskip
      When processing multiple examples at once, average the losses across all
      the examples, whether using MAE, MSE, or RMSE.
    \end{tdmargin}
  \end{withmargin}
\end{frame}

\begin{frame}{Choosing a loss: how should outliers count?}
  \begin{withmargin}
    \begin{tdmain}
      An outlier is a value far from the typical range. Cars are normally
      between 2000 and 5000 pounds and get between 8 to 50 miles per gallon. An
      8{,}000-pound car is an outlier.
      \vskip0.8em
      \textbf{MSE.} The model is closer to the outliers but further away from
      most of the other data points. L2 loss incurs a much higher penalty for an
      outlier than L1 loss.
      \vskip0.4em
      \textbf{MAE.} The model is further away from the outliers but closer to
      most of the other data points.
    \end{tdmain}
    \begin{tdmargin}
      {\color{tdfg}\textbf{Choose MSE}} if you want to heavily penalize large
      errors, or if you believe the outliers are important and indicative of
      true data variance the model should account for.
      \par\medskip
      {\color{tdfg}\textbf{Choose MAE}} if your dataset has significant outliers
      you don't want to overly influence the model, or if you prefer a loss that
      is directly interpretable as average error magnitude.
    \end{tdmargin}
  \end{withmargin}
  \keyterms{mean absolute error (MAE); mean squared error (MSE); L1; L2; loss;
    outlier; prediction}
\end{frame}

\begin{frame}{Gradient descent}
  \begin{withmargin}
    \begin{tdmain}
      Gradient descent is a mathematical technique that iteratively finds the
      weights and bias that produce the model with the lowest loss. The model
      begins training with randomized weights and biases near zero, then repeats
      the following steps.
      \vskip0.6em
      \begin{enumerate}
        \item Calculate the loss with the current weight and bias.
        \item Determine the direction to move the weights and bias that reduces loss.
        \item Move the weight and bias values a small amount in that direction.
        \item Return to step one and repeat until the model can't reduce the
          loss any further.
      \end{enumerate}
    \end{tdmain}
    \begin{tdmargin}
      The \gterm{small amount} in step three is the learning rate, covered two
      slides on.
      \par\medskip
      When a model \textbf{converges}, additional iterations don't reduce loss
      more, because gradient descent has found the weights and bias that nearly
      minimize the loss.
    \end{tdmargin}
  \end{withmargin}
\end{frame}

\begin{frame}{The math behind one step}
  \small
  Writing the prediction as $f_{w,b}(x) = (w \cdot x) + b$ and the actual value
  as $y$, MSE over $M$ examples is $\frac{1}{M}\sum_{i=1}^{M}(f_{w,b}(x_{(i)}) - y_{(i)})^2$.
  Taking the derivative with respect to the weight and the bias:
  \begin{align*}
    \frac{\partial}{\partial w}\frac{1}{M}\sum_{i=1}^{M}(f_{w,b}(x_{(i)}) - y_{(i)})^2
      &= \frac{1}{M}\sum_{i=1}^{M}(f_{w,b}(x_{(i)}) - y_{(i)}) \cdot 2x_{(i)}\\
    \frac{\partial}{\partial b}\frac{1}{M}\sum_{i=1}^{M}(f_{w,b}(x_{(i)}) - y_{(i)})^2
      &= \frac{1}{M}\sum_{i=1}^{M}(f_{w,b}(x_{(i)}) - y_{(i)}) \cdot 2\\[0.5em]
    \text{new weight} &= \text{old weight} - (\gterm{\text{small amount}} \cdot \text{weight slope})\\
    \text{new bias} &= \text{old bias} - (\gterm{\text{small amount}} \cdot \text{bias slope})
  \end{align*}
  \footnotesize
  Solving with weight and bias both zero on the seven-example car dataset gives
  a weight slope of $-119.7$ and a bias slope of $-34.3$.
\end{frame}

\begin{frame}{Six iterations on the car dataset}
  \begin{withmargin}
    \begin{tdmain}
      \footnotesize
      \begin{tabular}{@{}rrrr@{}}
        \toprule
        \textbf{Iteration} & \textbf{Weight} & \textbf{Bias} & \textbf{Loss (MSE)}\\
        \midrule
        1 & 0    & 0    & 303.71\\
        2 & 1.20 & 0.34 & 170.84\\
        3 & 2.05 & 0.59 & 103.17\\
        4 & 2.66 & 0.78 & 68.70\\
        5 & 3.09 & 0.91 & 51.13\\
        6 & 3.40 & 1.01 & \gterm{42.17}\\
        \bottomrule
      \end{tabular}
    \end{tdmain}
    \begin{tdmargin}
      Starting from weight 0 and bias 0, MSE is 303.71. With a small amount of
      0.01, the first update gives weight 1.2 and bias 0.34.
      \par\medskip
      In practice a model trains until it converges. Past convergence, loss
      fluctuates in small amounts as the model updates parameters around their
      lowest values, so train until the loss has stabilized.
    \end{tdmargin}
  \end{withmargin}
\end{frame}

\begin{frame}{Loss curves and convergence}
  \begin{withmargin}
    \begin{tdmain}
      A \gterm{loss curve} shows how the loss changes as the model trains. Loss
      is on the y-axis, iterations on the x-axis. Loss dramatically decreases
      during the first few iterations, then gradually decreases before flattening
      out around the 1{,}000th iteration mark.
      \vskip0.8em
      At around the second iteration the model would not be good at making
      predictions because of the high amount of loss. At around the 400th, a
      better model. At around the 1{,}000th, the model has converged.
    \end{tdmain}
    \begin{tdmargin}
      In the course's snapshots, blue loss lines run from the model to the data
      points and show the amount of loss. The longer the lines, the more loss
      there is.
    \end{tdmargin}
  \end{withmargin}
\end{frame}

\begin{frame}{Convergence and convex functions}
  \begin{withmargin}
    \begin{tdmain}
      The loss functions for linear models always produce a \gterm{convex}
      surface. As a result, when a linear regression model converges, we know
      the model has found the weights and bias that produce the lowest loss.
      \vskip0.8em
      If we graphed the weights and bias points during gradient descent, the
      points would look like a ball rolling down a hill, finally stopping at the
      point where there's no more downward slope.
      \vskip0.6em
      \footnotesize
      On a hypothetical miles-per-gallon dataset, a weight of $-5.44$ and bias of
      $35.94$ produce the lowest loss, at $5.54$.
    \end{tdmain}
    \begin{tdmargin}
      The model almost never finds the exact weight and bias values that
      minimize loss, but instead finds values very close to them.
      \par\medskip
      The minimum loss is typically greater than 0. A loss of 0 would mean the
      model fits every data point exactly, which is usually a sign of
      overfitting.
    \end{tdmargin}
  \end{withmargin}
  \keyterms{convergence; convex function; gradient descent; iteration; loss curve}
\end{frame}

\begin{frame}{Hyperparameters versus parameters}
  \vfill
  \begin{takeaway}
    Hyperparameters are values that \alertbf{you} control.\\
    Parameters are values that the \alertbf{model} calculates during training.
  \end{takeaway}
  \vfill
  {\small Three common hyperparameters: learning rate, batch size, and epochs.
  The weights and bias are parameters, part of the model itself.}
\end{frame}

\begin{frame}{Learning rate}
  \begin{withmargin}
    \begin{tdmain}
      Learning rate is a floating point number you set that influences how
      quickly the model converges. The model multiplies the gradient by the
      learning rate to determine the parameters for the next iteration.
      \vskip0.8em
      \begin{itemize}
        \item Too low: the model can take a very long time to converge, making
          only minor improvements after each iteration.
        \item Too high: the model never converges. It bounces around the weights
          and bias that minimize the loss, and loss can even increase at later
          iterations.
      \end{itemize}
      \vskip0.5em
      \footnotesize
      The \gterm{ideal learning rate is problem-dependent}. Each model and
      dataset has its own.
    \end{tdmain}
    \begin{tdmargin}
      The difference between old and new parameters is proportional to the slope
      of the loss function. If the slope is large, the model takes a large step.
      \par\medskip
      If the gradient's magnitude is 2.5 and the learning rate is 0.01, the
      model changes the parameter by 0.025.
    \end{tdmargin}
  \end{withmargin}
\end{frame}

\begin{frame}{Batch size}
  \begin{withmargin}
    \begin{tdmain}
      Batch size is the number of examples the model processes before updating
      its weights and bias. When a dataset contains hundreds of thousands or
      millions of examples, using the full batch isn't practical.
      \vskip0.7em
      \textbf{Stochastic gradient descent (SGD)} uses a single example per
      iteration. Given enough iterations SGD works, but is very noisy.
      \gterm{Noise} refers to variations during training that cause the loss to
      increase rather than decrease during an iteration.
      \vskip0.5em
      \textbf{Mini-batch SGD} is the compromise: for $N$ data points, batch size
      is any number greater than 1 and less than $N$. The model chooses the
      examples at random, averages their gradients, and updates once per
      iteration.
    \end{tdmain}
    \begin{tdmargin}
      Small batch sizes behave like SGD; larger batch sizes behave like
      full-batch gradient descent.
      \par\medskip
      A certain amount of noise can be a good thing. Noise can help a model
      generalize better and find the optimal weights and bias in a neural
      network.
      \par\medskip
      By averaging more gradients together, larger batch sizes help reduce the
      negative effects of outliers in the data.
    \end{tdmargin}
  \end{withmargin}
\end{frame}

\begin{frame}{Epochs}
  \begin{withmargin}
    \begin{tdmain}
      An \gterm{epoch} means the model has processed every example in the
      training set once. Given a training set with 1{,}000 examples and a
      mini-batch size of 100, it takes 10 iterations to complete one epoch.
      \vskip0.7em
      \footnotesize
      \begin{tabular}{@{}p{0.30\linewidth}p{0.64\linewidth}@{}}
        \toprule
        \textbf{Batch type} & \textbf{When updates occur, for 1{,}000 examples and 20 epochs}\\
        \midrule
        Full batch & After all examples: 20 updates, once per epoch\\[0.3em]
        Stochastic gradient descent & After a single example: 20{,}000 updates\\[0.3em]
        Mini-batch SGD (batch 100) & After each batch: 200 updates\\
        \bottomrule
      \end{tabular}
    \end{tdmain}
    \begin{tdmargin}
      The number of epochs is a hyperparameter you set before training. In many
      cases you'll need to experiment with how many epochs it takes to converge.
      \par\medskip
      In general, more epochs produces a better model, but also takes more time
      to train.
      \par\medskip
      Doubling the learning rate can slow down training, by making the weights
      bounce around.
    \end{tdmargin}
  \end{withmargin}
  \keyterms{batch size; epoch; generalize; hyperparameter; iteration; learning
    rate; mini-batch; mini-batch stochastic gradient descent; neural network;
    parameter; stochastic gradient descent}
\end{frame}

% ======================================================================
\sectionsubtitle{Predicting the probability of a given outcome.}
\section[Logistic regression]{Logistic Regression}

\begin{frame}{\modulehead{Logistic regression}{35 minutes}}
  \begin{withmargin}
    \begin{tdmain}
      In the linear regression module you explored how to construct a model to
      make continuous numerical predictions, such as the fuel efficiency of a
      car. But what if you want to build a model to answer questions like
      \gterm{Will it rain today?} or \gterm{Is this email spam?}
      \vskip0.9em
      Logistic regression is designed to predict the probability of a given
      outcome. Practically speaking, you can use the returned probability
      \textbf{as is}, or converted to a binary category such as \texttt{True} or
      \texttt{False}, \texttt{Spam} or \texttt{Not Spam}.
    \end{tdmain}
    \begin{tdmargin}
      If a spam-prediction model takes an email as input and outputs a value of
      \texttt{0.932}, this implies a \texttt{93.2\%} probability that the email
      is spam.
      \par\medskip
      This module uses the output as is. The classification module covers
      converting it into a binary category.
    \end{tdmargin}
  \end{withmargin}
\end{frame}

\begin{frame}{The sigmoid function}
  \begin{withmargin}
    \begin{tdmain}
      \[ f(x) = \frac{1}{1 + e^{-\gterm{x}}} \]
      \vskip0.5em
      \footnotesize
      As the input $x$ increases, the output approaches but never reaches 1.
      As the input decreases, the output approaches but never reaches 0. No
      matter how large or small the input, the output is always greater than 0
      and less than 1.
      \vskip0.5em
      \begin{tabular}{@{}rrrrrrrrr@{}}
        \toprule
        Input & $-7$ & $-3$ & $-1$ & 0 & 1 & 3 & 5 & 7\\
        \midrule
        Output & 0.001 & 0.047 & 0.269 & 0.50 & 0.731 & 0.952 & 0.993 & 0.999\\
        \bottomrule
      \end{tabular}
    \end{tdmain}
    \begin{tdmargin}
      There is a family of functions called logistic functions whose output is
      always between 0 and 1. The standard logistic function, also known as the
      \textbf{sigmoid} function, is the one above. Sigmoid means s-shaped.
      \par\medskip
      $e$ is Euler's number, a mathematical constant $\approx 2.71828$.
    \end{tdmargin}
  \end{withmargin}
\end{frame}

\begin{frame}{From linear output to probability}
  \begin{withmargin}
    \begin{tdmain}
      The linear component of a logistic regression model is
      \[ \gterm{z} = b + w_1x_1 + w_2x_2 + \ldots + w_Nx_N \]
      and $z$ is then passed to the sigmoid function, yielding a probability
      between 0 and 1:
      \[ y' = \frac{1}{1 + e^{-\gterm{z}}} \]
      \vskip0.3em
      \footnotesize
      A linear equation can output very big or very small values of $z$, but the
      output $y'$ is always between 0 and 1, exclusive. A $z$ value of $-10$
      maps to a $y'$ value of 0.00004.
    \end{tdmain}
    \begin{tdmargin}
      $z$ is called the \textbf{log-odds}, because solving the sigmoid for $z$
      gives
      \par\smallskip
      $z = \ln\left(\frac{y}{1-y}\right)$
      \par\smallskip
      the natural logarithm of the ratio of the probabilities of the two
      possible outcomes, $y$ and $1-y$.
      \par\medskip
      The sigmoid bends the straight line into an s-shape.
    \end{tdmargin}
  \end{withmargin}
\end{frame}

\begin{frame}{Worked example}
  \small
  A logistic regression model with three features has bias $b = 1$ and weights
  $w_1 = 2$, $w_2 = -1$, $w_3 = 5$. Given inputs $x_1 = 0$, $x_2 = 10$,
  $x_3 = 2$:
  \begin{align*}
    z &= 1 + 2x_1 - x_2 + 5x_3\\
      &= 1 + (2)(0) - (10) + (5)(2)\\
      &= \gterm{1}\\[0.6em]
    y &= \frac{1}{1 + e^{-z}} = \frac{1}{1 + e^{-1}} = \frac{1}{1 + 0.367}
       = \frac{1}{1.367} = \gterm{0.731}
  \end{align*}
  \footnotesize
  Remember, the output of the sigmoid function will always be greater than 0 and
  less than 1.
  \keyterms{binary classification; log odds; logistic regression; sigmoid function}
\end{frame}

\begin{frame}{Why not squared loss}
  \begin{withmargin}
    \begin{tdmain}
      Squared loss works well for a linear model where the rate of change of the
      output values is constant. Given $y' = b + 3x_1$, each time you increment
      $x_1$ by 1, $y'$ increases by 3.
      \vskip0.6em
      The rate of change of a logistic regression model is \alertbf{not
      constant}. When $z$ is close to 0, small increases in $z$ result in much
      larger changes to $y$ than when $z$ is a large positive or negative number.
      \vskip0.5em
      \footnotesize
      \begin{tabular}{@{}rrrrrr@{}}
        \toprule
        Input & 5 & 7 & 8 & 9 & 10\\
        \midrule
        Logistic output & 0.993 & 0.999 & 0.9997 & 0.9999 & 0.99998\\
        Digits of precision & 3 & 3 & 4 & 4 & 5\\
        \bottomrule
      \end{tabular}
    \end{tdmain}
    \begin{tdmargin}
      If you used squared loss for the sigmoid, then as the output got closer
      and closer to \texttt{0} and \texttt{1}, you would need more memory to
      preserve the precision needed to track these values.
    \end{tdmargin}
  \end{withmargin}
\end{frame}

\begin{frame}{Log Loss}
  \begin{withmargin}
    \begin{tdmain}
      \[ \text{Log Loss} = -\frac{1}{N}\sum_{i=1}^{N}
         \left[\gterm{y_i}\log(y_i') + (1 - \gterm{y_i})\log(1 - y_i')\right] \]
      \vskip0.5em
      \footnotesize
      \begin{itemize}
        \item $N$ is the number of labeled examples in the dataset.
        \item $i$ is the index of an example in the dataset.
        \item $\gterm{y_i}$ is the label for the $i$th example. Since this is
          logistic regression, it must either be 0 or 1.
        \item $y_i'$ is the model's prediction for the $i$th example, somewhere
          between 0 and 1.
      \end{itemize}
    \end{tdmain}
    \begin{tdmargin}
      The Log Loss equation returns the logarithm of the magnitude of the
      change, rather than just the distance from data to prediction.
      \par\medskip
      This form calculates the \textbf{mean} Log Loss across all points. Using
      mean rather than total Log Loss is desirable in practice, because it
      decouples tuning of the batch size and the learning rate.
    \end{tdmargin}
  \end{withmargin}
\end{frame}

\begin{frame}{Regularization in logistic regression}
  \begin{withmargin}
    \begin{tdmain}
      Regularization, a mechanism for penalizing model complexity during
      training, is extremely important in logistic regression modeling. Without
      it, the \alertbf{asymptotic} nature of logistic regression would keep
      driving loss towards 0 in cases where the model has a large number of
      features.
      \vskip0.8em
      Most logistic regression models use one of two strategies to decrease
      model complexity:
      \vskip0.4em
      \begin{itemize}
        \item L2 regularization.
        \item Early stopping: limiting the number of training steps to halt
          training while loss is still decreasing.
      \end{itemize}
    \end{tdmain}
    \begin{tdmargin}
      Logistic regression models are trained using the same process as linear
      regression models, with two distinctions: Log Loss instead of squared
      loss, and regularization is critical.
    \end{tdmargin}
  \end{withmargin}
  \keyterms{gradient descent; linear regression; Log Loss; logistic regression;
    loss function; overfitting; regularization; squared loss}
\end{frame}

% ======================================================================
\sectionsubtitle{From a probability to a category, and how to judge it.}
\section[Classification]{Classification}

\begin{frame}{\modulehead{Classification}{70 minutes}}
  \begin{withmargin}
    \begin{tdmain}
      In logistic regression you learned to use the sigmoid function to convert
      raw model output to a value between 0 and 1, for example predicting that a
      given email has a 75\% chance of being spam. But what if your goal is not
      to output a probability but a \gterm{category}?
      \vskip0.9em
      Classification is the task of predicting which of a set of classes an
      example belongs to. This module converts a logistic regression model that
      predicts a probability into a binary classification model that predicts
      one of two classes, and shows how to choose and calculate metrics to
      evaluate the quality of its predictions.
    \end{tdmain}
    \begin{tdmargin}
      {\color{tdfg}\textbf{Objectives.}}
      \par\medskip
      Determine an appropriate threshold for a binary classification model.
      Calculate and choose appropriate metrics to evaluate one. Interpret ROC
      and AUC.
    \end{tdmargin}
  \end{withmargin}
\end{frame}

\begin{frame}{The classification threshold}
  \begin{withmargin}
    \begin{tdmain}
      To convert the model's raw numerical output into one of two categories,
      you choose a threshold probability, called a \gterm{classification
      threshold}. Examples with a probability above the threshold are assigned
      to the \textbf{positive class}, the class you are testing for. Examples
      with a lower probability are assigned to the \textbf{negative class}.
      \vskip0.8em
      While 0.5 might seem like an intuitive threshold, it's not a good idea if
      the cost of one type of wrong classification is greater than the other, or
      if the classes are imbalanced.
    \end{tdmain}
    \begin{tdmargin}
      Suppose the model scores one email at 0.99 and another at 0.51. With a
      threshold of 0.5, both are spam. With a threshold of 0.95, only the 0.99
      email is spam.
      \par\medskip
      If only 0.01\% of emails are spam, or if misfiling legitimate emails is
      worse than letting spam into the inbox, labelling anything at least 50\%
      likely as spam produces undesirable results.
    \end{tdmargin}
  \end{withmargin}
\end{frame}

\begin{frame}{The confusion matrix}
  \begin{withmargin}
    \begin{tdmain}
      \footnotesize
      \begin{tabular}{@{}p{0.20\linewidth}p{0.36\linewidth}p{0.36\linewidth}@{}}
        \toprule
         & \textbf{Actual positive} & \textbf{Actual negative}\\
        \midrule
        \textbf{Predicted positive}
          & True positive (TP): a spam email correctly classified as spam.
            Sent to the spam folder.
          & False positive (FP): a not-spam email misclassified as spam.
            Legitimate mail in the spam folder.\\[0.5em]
        \textbf{Predicted negative}
          & False negative (FN): a spam email misclassified as not-spam. Spam
            that reaches the inbox.
          & True negative (TN): a not-spam email correctly classified. Sent to
            the inbox.\\
        \bottomrule
      \end{tabular}
    \end{tdmain}
    \begin{tdmargin}
      Row totals give all predicted positives (TP + FP) and all predicted
      negatives (FN + TN), regardless of validity. Column totals give all real
      positives (TP + FN) and all real negatives (FP + TN), regardless of model
      classification.
      \par\medskip
      When the total of actual positives is not close to the total of actual
      negatives, the dataset is \gterm{imbalanced}.
    \end{tdmargin}
  \end{withmargin}
\end{frame}

\begin{frame}{What moving the threshold does}
  \begin{withmargin}
    \begin{tdmain}
      As the threshold \textbf{increases}, the model predicts fewer positives
      overall, both true and false, and more negatives overall, both true and
      false.
      \vskip0.9em
      A spam classifier with a threshold of \gterm{0.9999} will only label an
      email as spam if it considers the classification at least 99.99\% likely.
      It is highly unlikely to mislabel a legitimate email, but also likely to
      miss actual spam. At a very high threshold, almost all emails, both spam
      and not-spam, will be classified as not-spam.
    \end{tdmain}
    \begin{tdmargin}
      The course's threshold widget carries three toy datasets:
      \par\medskip
      \textbf{Separated}, where positive and negative examples are generally
      well differentiated.
      \par\smallskip
      \textbf{Unseparated}, where many positives score lower than negatives and
      vice versa.
      \par\smallskip
      \textbf{Imbalanced}, containing only a few examples of the positive class.
    \end{tdmargin}
  \end{withmargin}
  \keyterms{binary classification; class-imbalanced dataset; classification
    threshold; confusion matrix; ground truth; negative class; positive class}
\end{frame}

\begin{frame}{Accuracy, and when it lies}
  \begin{withmargin}
    \begin{tdmain}
      \[ \text{Accuracy} = \frac{\text{correct classifications}}{\text{total classifications}}
         = \frac{TP + TN}{TP + TN + FP + FN} \]
      \vskip0.5em
      Given a balanced dataset, accuracy can serve as a coarse-grained measure
      of model quality. For heavily imbalanced datasets, where one class appears
      very rarely, say 1\% of the time, a model that predicts negative 100\% of
      the time would score \alertbf{99\% on accuracy, despite being useless}.
    \end{tdmain}
    \begin{tdmargin}
      All of the metrics in this section are calculated at a single fixed
      threshold, and change when the threshold changes. Very often, the user
      tunes the threshold to optimize one of them.
      \par\medskip
      In ML, words like recall, precision, and accuracy have mathematical
      definitions that may differ from, or be more specific than, more commonly
      used meanings of the word.
    \end{tdmargin}
  \end{withmargin}
\end{frame}

\begin{frame}{Recall, false positive rate, precision}
  \begin{withmargin}
    \begin{tdmain}
      \begin{align*}
        \text{Recall (TPR)} &= \frac{\text{correctly classified actual positives}}{\text{all actual positives}} = \frac{TP}{TP + FN}\\[0.4em]
        \text{FPR} &= \frac{\text{incorrectly classified actual negatives}}{\text{all actual negatives}} = \frac{FP}{FP + TN}\\[0.4em]
        \text{Precision} &= \frac{\text{correctly classified actual positives}}{\text{everything classified as positive}} = \frac{TP}{TP + FP}
      \end{align*}
      \vskip0.3em
      \footnotesize
      Precision improves as false positives decrease; recall improves as false
      negatives decrease. Raising the threshold decreases false positives and
      increases false negatives, so precision and recall often show an
      \gterm{inverse relationship}.
    \end{tdmain}
    \begin{tdmargin}
      Another name for recall is \textbf{probability of detection}: what
      fraction of spam emails are detected by this model?
      \par\medskip
      FPR is also the \textbf{probability of false alarm}: what fraction of
      legitimate emails were incorrectly classified as spam?
      \par\medskip
      Precision: what fraction of emails classified as spam were actually spam?
    \end{tdmargin}
  \end{withmargin}
\end{frame}

\begin{frame}{Which metric, and when}
  \begin{withmargin}
    \begin{tdmain}
      \footnotesize
      \begin{tabular}{@{}p{0.26\linewidth}p{0.68\linewidth}@{}}
        \toprule
        \textbf{Metric} & \textbf{Guidance}\\
        \midrule
        Accuracy & Use as a rough indicator of training progress for balanced
          datasets. For model performance, use only in combination with other
          metrics. Avoid for imbalanced datasets.\\[0.4em]
        Recall (TPR) & Use when false negatives are more expensive than false
          positives.\\[0.4em]
        False positive rate & Use when false positives are more expensive than
          false negatives.\\[0.4em]
        Precision & Use when it's very important for positive predictions to be
          accurate.\\
        \bottomrule
      \end{tabular}
    \end{tdmain}
    \begin{tdmargin}
      \gterm{NaN}, not a number, appears when dividing by 0. When TP and FP are
      both 0, precision has 0 in the denominator. A model that never predicts
      positive would have 0 TPs and 0 FPs, so its precision is NaN. NaN can look
      like perfect performance but can also come from a useless model.
    \end{tdmargin}
  \end{withmargin}
\end{frame}

\begin{frame}{F1 score}
  \begin{withmargin}
    \begin{tdmain}
      \[ \text{F1} = 2 \cdot \frac{\text{precision} \cdot \text{recall}}{\text{precision} + \text{recall}}
         = \frac{2TP}{2TP + FP + FN} \]
      \vskip0.6em
      The F1 score is the \gterm{harmonic mean} of precision and recall. This
      metric balances the importance of both, and is preferable to accuracy for
      class-imbalanced datasets.
    \end{tdmain}
    \begin{tdmargin}
      When precision and recall both have perfect scores of 1.0, F1 will also be
      1.0.
      \par\medskip
      When precision and recall are close in value, F1 will be close to their
      value. When they are far apart, F1 will be similar to whichever metric is
      worse.
    \end{tdmargin}
  \end{withmargin}
  \keyterms{accuracy; false positive rate (FPR); precision; recall}
\end{frame}

\begin{frame}{ROC and AUC}
  \begin{withmargin}
    \begin{tdmain}
      The ROC curve is a visual representation of model performance across
      \alertbf{all} thresholds. It is drawn by calculating TPR and FPR at every
      possible threshold, then graphing TPR over FPR.
      \vskip0.8em
      The area under the ROC curve (AUC) represents the probability that the
      model, if given a randomly chosen positive and negative example, will rank
      the positive higher than the negative.
      \vskip0.6em
      \footnotesize
      A perfect model has an AUC of 1.0. A model that does exactly as well as
      coin flips has a diagonal ROC from (0,0) to (1,1) and an AUC of 0.5.
    \end{tdmain}
    \begin{tdmargin}
      The long version of the name, \textbf{receiver operating characteristic},
      is a holdover from WWII radar detection.
      \par\medskip
      A spam classifier with AUC of 1.0 always assigns a random spam email a
      higher probability of being spam than a random legitimate email. With AUC
      of 0.5, it does so only half the time.
    \end{tdmargin}
  \end{withmargin}
\end{frame}

\begin{frame}{Choosing a model, and a threshold, with AUC}
  \begin{withmargin}
    \begin{tdmain}
      AUC is a useful measure for comparing two models, as long as the dataset
      is roughly balanced. The model with greater area under the curve is
      generally the better one. A model with AUC below 0.5 performs
      \gterm{worse than chance}.
      \vskip0.8em
      The points on a ROC curve closest to (0,1) represent the best-performing
      thresholds. If false alarms are highly costly, choose a threshold with a
      lower FPR even if TPR is reduced. If false positives are cheap and missed
      true positives highly costly, choose the threshold that maximizes TPR.
    \end{tdmain}
    \begin{tdmargin}
      {\color{tdfg}\textbf{Fixing a worse-than-chance model.}}
      \par\medskip
      Reverse the predictions, so 1 becomes 0 and 0 becomes 1. If a binary
      classifier reliably puts examples in the wrong classes more often than
      chance, switching the class label immediately makes its predictions better
      than chance without retraining.
      \par\medskip
      When the dataset is imbalanced, precision-recall curves may offer a better
      comparative visualization.
    \end{tdmargin}
  \end{withmargin}
  \keyterms{area under the ROC curve (AUC); ROC curve}
\end{frame}

\begin{frame}{Prediction bias}
  \begin{withmargin}
    \begin{tdmain}
      Prediction bias is the difference between the mean of a model's
      predictions and the mean of ground-truth labels in the data. A model
      trained on a dataset where 5\% of the emails are spam should predict, on
      average, that \gterm{5\%} of the emails it classifies are spam.
      \vskip0.8em
      If the model instead predicts 50\% of the time that an email is spam, then
      something is wrong with the training dataset, the new dataset the model is
      applied to, or the model itself.
    \end{tdmain}
    \begin{tdmargin}
      {\color{tdfg}\textbf{Causes.}}
      \par\medskip
      Biases or noise in the data, including biased sampling for the training
      set.
      \par\smallskip
      Too-strong regularization, meaning the model was oversimplified and lost
      necessary complexity.
      \par\smallskip
      Bugs in the model training pipeline.
      \par\smallskip
      The set of features provided being insufficient for the task.
    \end{tdmargin}
  \end{withmargin}
\end{frame}

\begin{frame}{Multi-class classification}
  \begin{withmargin}
    \begin{tdmain}
      Multi-class classification is an extension of binary classification to
      more than two classes. If each example can only be assigned to one class,
      the problem can be handled as a series of binary classification problems,
      where one class contains one of the multiple classes and the other
      contains all the rest put together.
      \vskip0.8em
      For labels A, B, and C: first build a classifier that categorizes examples
      as A+B versus C. Then build a second that reclassifies the A+B examples as
      A versus B.
      \vskip0.5em
      \footnotesize
      If class membership isn't exclusive, so that an example can be assigned to
      multiple classes, this is a \gterm{multi-label} classification problem.
    \end{tdmain}
    \begin{tdmargin}
      An example of a multi-class problem is a handwriting classifier that takes
      an image of a handwritten digit and decides which digit, 0 to 9, is
      represented.
      \par\medskip
      Multi-class classification is explored more deeply in the neural networks
      module.
    \end{tdmargin}
  \end{withmargin}
  \keyterms{ground truth; prediction; prediction bias; binary classification}
\end{frame}

% ======================================================================
\sectionsubtitle{The model ingests a feature vector, not the dataset.}
\section[Numerical data]{Working with Numerical Data}

\begin{frame}{\modulehead{Working with numerical data}{85 minutes}}
  \begin{withmargin}
    \begin{tdmain}
      This unit focuses on numerical data, meaning integers or floating-point
      values that behave like numbers. That is, they are additive, countable,
      ordered, and so on. Examples: temperature, weight, the number of deer
      wintering in a nature preserve.
      \vskip0.9em
      In contrast, US postal codes, despite being five-digit or nine-digit
      numbers, don't behave like numbers. Postal code 40004 is not twice the
      quantity of postal code 20002. These numbers represent categories,
      specifically geographic areas, and are \gterm{categorical data}.
    \end{tdmain}
    \begin{tdmargin}
      {\color{tdfg}\textbf{Three units on data.}}
      \par\medskip
      Data is so important that this course devotes three entire units to the
      topic: working with numerical data, working with categorical data, and
      datasets, generalization, and overfitting.
    \end{tdmargin}
  \end{withmargin}
\end{frame}

\begin{frame}{Feature vectors}
  \begin{withmargin}
    \begin{tdmain}
      A model does not act directly on the rows of a dataset. The model ingests
      an array of floating-point values called a \gterm{feature vector}. You can
      think of a feature vector as the floating-point values comprising one
      example.
      \vskip0.8em
      Feature vectors seldom use the dataset's raw values. You must typically
      process the dataset's values into representations that your model can
      better learn from. Would a model produce better predictions by training
      from the actual values rather than altered values? Surprisingly, the
      answer is \alertbf{no}.
    \end{tdmain}
    \begin{tdmargin}
      Determining the best way to represent raw dataset values as trainable
      values in the feature vector is called \textbf{feature engineering}, and
      it is a vital part of machine learning.
      \par\medskip
      The most common techniques are \textbf{normalization}, converting
      numerical values into a standard range, and \textbf{binning}, converting
      them into buckets of ranges.
    \end{tdmargin}
  \end{withmargin}
  \keyterms{binning; bucketing; feature engineering; feature vector;
    normalization; preprocessing}
\end{frame}

\begin{frame}{First steps: look, then measure}
  \begin{withmargin}
    \begin{tdmain}
      Before creating feature vectors, study numerical data in two ways.
      \vskip0.6em
      \textbf{Visualize.} Graphs help you find anomalies or patterns hiding in
      the data. Look at your data graphically, either as scatter plots or
      histograms, not only at the beginning of the pipeline but throughout data
      transformations. Visualizations help you continually check your
      assumptions.
      \vskip0.5em
      \textbf{Evaluate statistically.} Gather mean and median, standard
      deviation, and the values at the quartile divisions: the 0th, 25th, 50th,
      75th, and 100th percentiles.
    \end{tdmain}
    \begin{tdmargin}
      The 0th percentile is the minimum value of the column; the 100th is the
      maximum. The 50th percentile is the median.
      \par\medskip
      Certain visualization tools are optimized for certain data formats. A tool
      that helps you evaluate protocol buffers may or may not help you evaluate
      CSV data.
    \end{tdmargin}
  \end{withmargin}
\end{frame}

\begin{frame}{Finding outliers}
  \begin{withmargin}
    \begin{tdmain}
      An outlier is a value distant from most other values in a feature or
      label. Outliers often cause problems in model training, so finding
      outliers is important.
      \vskip0.7em
      When the delta between the 0th and 25th percentiles differs significantly
      from the delta between the 75th and 100th percentiles, the dataset
      \gterm{probably contains outliers}.
      \vskip0.7em
      \footnotesize
      \begin{itemize}
        \item The outlier is due to a mistake. An experimenter entered an extra
          zero, or an instrument malfunctioned. Generally, delete these examples.
        \item The outlier is legitimate. Ask: will your trained model ultimately
          need to infer good predictions on these outliers?
      \end{itemize}
    \end{tdmain}
    \begin{tdmargin}
      If yes, keep these outliers in your training set. Outliers in certain
      features sometimes mirror outliers in the label, so they could help your
      model make better predictions. Be careful: extreme outliers can still hurt
      your model.
      \par\medskip
      If no, delete the outliers or apply more invasive feature engineering
      techniques, such as clipping.
      \par\medskip
      Don't over-rely on basic statistics. Anomalies can also hide in seemingly
      well-balanced data.
    \end{tdmargin}
  \end{withmargin}
  \keyterms{clipping; outliers}
\end{frame}

\begin{frame}{Why normalize}
  \begin{withmargin}
    \begin{tdmain}
      The goal of normalization is to transform features to be on a similar
      scale. Feature \texttt{X} spans 154 to 24{,}917{,}482; feature \texttt{Y}
      spans 5 to 22. Normalization might manipulate them so they span a similar
      range, perhaps 0 to 1.
      \vskip0.6em
      \footnotesize
      \begin{itemize}
        \item Models converge more quickly. When features have different ranges,
          gradient descent can bounce and slow convergence.
        \item Models infer better predictions.
        \item Avoids the \gterm{NaN trap}. When a value exceeds the
          floating-point precision limit, the system sets it to \texttt{NaN},
          and other numbers eventually become NaN too.
        \item The model learns appropriate weights for each feature. Without
          scaling, the model pays too much attention to features with wide
          ranges.
      \end{itemize}
    \end{tdmain}
    \begin{tdmargin}
      {\color{tdfg}\textbf{Warning.}} If you normalize a feature during
      training, you must also normalize that feature when making predictions.
      \par\medskip
      Advanced optimizers like Adagrad and Adam protect against the convergence
      problem by changing the effective learning rate over time.
    \end{tdmargin}
  \end{withmargin}
\end{frame}

\begin{frame}{How much disparity matters}
  \begin{withmargin}
    \begin{tdmain}
      Feature \texttt{A} spans $-0.5$ to $+0.5$; feature \texttt{B} spans $-5.0$
      to $+5.0$. Both spans are narrow, but \texttt{B}'s is ten times wider.
      Therefore, at the start of training, the model assumes \texttt{B} is ten
      times more important than \texttt{A}; training takes longer than it
      should; the resulting model may be suboptimal. The overall damage is
      relatively small, but still normalize both to the same scale.
      \vskip0.8em
      Now feature \texttt{C} spans $-1$ to $+1$ and feature \texttt{D} spans
      $+5000$ to $+1{,}000{,}000{,}000$. Without normalization the model will
      likely be suboptimal, and training will take much longer to converge or
      \alertbf{fail to converge entirely}.
    \end{tdmain}
    \begin{tdmargin}
      Normalize numeric features covering distinctly different ranges, for
      example age and income.
      \par\medskip
      Also normalize a single numeric feature that covers a wide range, such as
      \texttt{city population}.
    \end{tdmargin}
  \end{withmargin}
\end{frame}

\begin{frame}[fragile]{Linear scaling}
  \begin{withmargin}
    \begin{tdmain}
      \[ x' = \frac{x - \gterm{x_{min}}}{\gterm{x_{max}} - \gterm{x_{min}}} \]
      \vskip0.4em
      Linear scaling converts floating-point values from their natural range
      into a standard range, usually 0 to 1 or $-1$ to $+1$. For a feature
      \texttt{quantity} whose natural range spans 100 to 900, a value of 300
      normalizes as follows.
      \vskip0.3em
      \begin{tdcode}
\begin{Verbatim}[fontsize=\scriptsize]
x' = (300 - 100) / (900 - 100)
x' = 200 / 800
x' = 0.25
\end{Verbatim}
      \end{tdcode}
    \end{tdmain}
    \begin{tdmargin}
      Good when the lower and upper bounds don't change much over time, the
      feature contains few or no extreme outliers, and the feature is
      approximately uniformly distributed across its range.
      \par\medskip
      Human \texttt{age} qualifies: bounds are roughly 0 to 100, only about
      0.3\% of the population is over 100, and a large dataset should contain
      sufficient examples of all ages.
      \par\medskip
      \texttt{net\_worth} does not: many outliers, values not uniformly
      distributed.
    \end{tdmargin}
  \end{withmargin}
\end{frame}

\begin{frame}[fragile]{Z-score scaling}
  \begin{withmargin}
    \begin{tdmain}
      \[ x' = \frac{x - \gterm{\mu}}{\gterm{\sigma}} \]
      \vskip0.4em
      A Z-score is the number of standard deviations a value is from the mean. A
      value 2 standard deviations greater than the mean has a Z-score of $+2.0$.
      \vskip0.3em
      \begin{tdcode}
\begin{Verbatim}[fontsize=\scriptsize]
mean = 100, standard deviation = 20
original value = 130

Z-score = (130 - 100) / 20
Z-score = 30 / 20
Z-score = +1.5
\end{Verbatim}
      \end{tdcode}
    \end{tdmain}
    \begin{tdmargin}
      Z-score is a good choice when the data follows a normal distribution or
      something like one. It is typically a better choice than linear scaling,
      because most real-world features do not meet all the criteria for linear
      scaling.
      \par\medskip
      Adult \texttt{height} over ten million women qualifies: it conforms to a
      normal distribution, and ten million examples implies enough outliers for
      the model to learn patterns on very high or very low Z-scores.
    \end{tdmargin}
  \end{withmargin}
\end{frame}

\begin{frame}{How rare is a Z-score of 4?}
  \begin{withmargin}
    \begin{tdmain}
      In a classic normal distribution:
      \vskip0.5em
      \begin{itemize}
        \item At least 68.27\% of data has a Z-score between $-1.0$ and $+1.0$.
        \item At least 95.45\% between $-2.0$ and $+2.0$.
        \item At least 99.73\% between $-3.0$ and $+3.0$.
        \item At least 99.994\% between $-4.0$ and $+4.0$.
      \end{itemize}
      \vskip0.6em
      \footnotesize
      A feature with one billion examples conforming to a classic normal
      distribution could still have as many as \gterm{60{,}000} examples outside
      the range $-4.0$ to $+4.0$.
    \end{tdmain}
    \begin{tdmargin}
      So are those points truly outliers? Outliers is a concept without a strict
      definition, so no one can say for sure.
      \par\medskip
      Some distributions are normal within the bulk of their range but contain
      extreme outliers. Nearly all points in \texttt{net\_worth} might fit within
      3 standard deviations, but a few could be hundreds away. Combine Z-score
      scaling with clipping.
    \end{tdmargin}
  \end{withmargin}
\end{frame}

\begin{frame}[fragile]{Log scaling}
  \begin{withmargin}
    \begin{tdmain}
      \[ x' = \gterm{\ln}(x) \]
      \vskip0.4em
      Log scaling is helpful when the data conforms to a \textbf{power law}
      distribution: low values of \texttt{X} have very high values of
      \texttt{Y}, and as \texttt{X} increases, \texttt{Y} quickly decreases.
      \vskip0.4em
      \begin{tdcode}
\begin{Verbatim}[fontsize=\scriptsize]
~4.6  = ln(100)
~13.8 = ln(1,000,000)
\end{Verbatim}
      \end{tdcode}
      \vskip0.3em
      \footnotesize
      The log of 1{,}000{,}000 is only about three times larger than the log of
      100.
    \end{tdmain}
    \begin{tdmargin}
      Movie ratings: a few movies have lots of user ratings, most have very few.
      \par\medskip
      Book sales: most published books sell one or two hundred copies, some sell
      thousands, only a few bestsellers sell more than a million.
      \par\medskip
      A linear model on raw sales would have to find something about a
      million-selling book cover that is 10{,}000 times more powerful than a
      100-selling cover. You could imagine it being about three times more
      powerful.
    \end{tdmargin}
  \end{withmargin}
\end{frame}

\begin{frame}{Clipping}
  \begin{withmargin}
    \begin{tdmain}
      Clipping caps the value of outliers to a specific maximum value. Clipping
      is a strange idea, and yet, it can be very effective.
      \vskip0.8em
      For a feature \texttt{roomsPerPerson}, over 99\% of values conform to a
      normal distribution with mean about 1.8 and standard deviation about 0.7,
      but a few outliers reach all the way out to 17. Capping at 4.0 does not
      mean the model ignores values greater than 4.0. It means all values that
      were greater than 4.0 \gterm{now become 4.0}. That explains the peculiar
      hill at 4.0, and the scaled feature set is now more useful than the
      original data.
    \end{tdmain}
    \begin{tdmargin}
      You can also clip after applying other normalization. If you use Z-score
      scaling but a few outliers have absolute values far greater than 3, clip
      Z-scores greater than 3 to exactly 3 and less than $-3$ to exactly $-3$.
      \par\medskip
      Clipping prevents your model from overindexing on unimportant data.
      However, some outliers are actually important, so clip values carefully.
    \end{tdmargin}
  \end{withmargin}
\end{frame}

\begin{frame}{The four techniques, by distribution shape}
  \begin{withmargin}
    \begin{tdmain}
      \footnotesize
      \begin{tabular}{@{}p{0.24\linewidth}p{0.30\linewidth}p{0.38\linewidth}@{}}
        \toprule
        \textbf{Technique} & \textbf{Formula} & \textbf{When to use}\\
        \midrule
        Linear scaling & $x' = \frac{x - x_{min}}{x_{max} - x_{min}}$
          & Mostly uniformly distributed across range. Flat-shaped.\\[0.5em]
        Z-score scaling & $x' = \frac{x - \mu}{\sigma}$
          & Normally distributed, peak close to mean. Bell-shaped.\\[0.5em]
        Log scaling & $x' = \log(x)$
          & Heavily skewed on at least one tail. Heavy-tail-shaped.\\[0.5em]
        Clipping & If $x > max$, set $x' = max$; if $x < min$, set $x' = min$
          & Contains extreme outliers.\\
        \bottomrule
      \end{tabular}
    \end{tdmain}
    \begin{tdmargin}
      The best normalization technique is one that works well in practice, so
      try new ideas if you think they'll work well on your feature distribution.
      \par\medskip
      \gterm{Mistakes get deleted, legitimate outliers get clipped.} For a data
      centre \texttt{temperature} feature that is normally 15 to 30, clip the
      legitimate values between 31 and 45, but delete the every-1000th values of
      1{,}000.
    \end{tdmargin}
  \end{withmargin}
  \keyterms{clipping; NaN trap; normalization; scaling; Z-score normalization}
\end{frame}

\begin{frame}{Binning}
  \begin{withmargin}
    \begin{tdmain}
      Binning, also called bucketing, groups numerical subranges into bins. In
      many cases binning turns numerical data into categorical data. A feature
      \texttt{X} spanning 15 to 425 might become five bins.
      \vskip0.5em
      \footnotesize
      \begin{tabular}{@{}rlr@{}}
        \toprule
        \textbf{Bin} & \textbf{Range} & \textbf{Feature vector}\\
        \midrule
        1 & 15--34   & [1.0, 0.0, 0.0, 0.0, 0.0]\\
        2 & 35--117  & [0.0, 1.0, 0.0, 0.0, 0.0]\\
        3 & 118--279 & [0.0, 0.0, 1.0, 0.0, 0.0]\\
        4 & 280--392 & [0.0, 0.0, 0.0, 1.0, 0.0]\\
        5 & 393--425 & [0.0, 0.0, 0.0, 0.0, 1.0]\\
        \bottomrule
      \end{tabular}
    \end{tdmain}
    \begin{tdmargin}
      A model trained on these bins will react no differently to \texttt{X}
      values of 17 and 29, since both are in bin 1.
      \par\medskip
      Even though \texttt{X} is a single column, binning causes a model to treat
      it as \gterm{five separate features}, and the model learns separate
      weights for each bin.
    \end{tdmargin}
  \end{withmargin}
\end{frame}

\begin{frame}{When binning beats scaling}
  \begin{withmargin}
    \begin{tdmain}
      Binning is a good alternative to scaling or clipping when the overall
      linear relationship between the feature and the label is weak or
      nonexistent, or when the feature values are clustered.
      \vskip0.8em
      Predicting the number of shoppers from outside temperature: the plot shows
      three clusters, at 4 to 11, 12 to 26, and 27 to 36 degrees. Represented as
      one raw feature, a temperature of 35.0 would have five times the influence
      of a temperature of 7.0, but the plot shows \alertbf{no linear
      relationship} at all.
    \end{tdmain}
    \begin{tdmargin}
      Binning can feel counterintuitive, given that the model treats the values
      37 and 115 identically. But when a feature appears more clumpy than
      linear, binning is a much better way to represent the data.
      \par\medskip
      More bins is usually a bad idea. A model can only learn the association
      between a bin and a label if there are enough examples in that bin, and
      you should minimize the number of features in a model.
    \end{tdmargin}
  \end{withmargin}
\end{frame}

\begin{frame}{Quantile bucketing}
  \begin{withmargin}
    \begin{tdmain}
      Quantile bucketing creates boundaries such that the number of examples in
      each bucket is exactly or nearly equal. Quantile bucketing mostly
      \gterm{hides the outliers}.
      \vskip0.8em
      With ten equally spaced 10{,}000-dollar car-price buckets, the bucket from
      0 to 10{,}000 contains dozens of examples but the bucket from 50{,}000 to
      60{,}000 contains only 5. The model has enough examples to train on the
      first but not the second. Quantile buckets instead give the first bucket
      0 to 4{,}000, the second 4{,}001 to 6{,}000, and the sixth 25{,}001 to
      60{,}000.
    \end{tdmain}
    \begin{tdmargin}
      Bucketing with equal intervals works for many data distributions. For
      skewed data, try quantile bucketing.
      \par\medskip
      Equal intervals give extra information space to the long tail while
      compacting the large torso into a single bucket. Quantile buckets give
      extra space to the large torso while compacting the long tail.
    \end{tdmargin}
  \end{withmargin}
  \keyterms{binning; bucketing; clipping; feature; feature engineering; feature
    vector; label; scaling}
\end{frame}

\begin{frame}{Scrubbing}
  \begin{withmargin}
    \begin{tdmain}
      Apple trees produce a mixture of great fruit and wormy messes. Yet the
      apples in high-end grocery stores display 100\% perfect fruit. As an ML
      engineer, you'll spend enormous amounts of your time tossing out bad
      examples and cleaning up the salvageable ones. \alertbf{Even a few bad
      apples can spoil a large dataset.}
      \vskip0.6em
      \footnotesize
      \begin{tabular}{@{}p{0.36\linewidth}p{0.56\linewidth}@{}}
        \toprule
        \textbf{Problem category} & \textbf{Example}\\
        \midrule
        Omitted values & A census taker fails to record a resident's age\\
        Duplicate examples & A server uploads the same logs twice\\
        Out-of-range feature values & A human accidentally types an extra digit\\
        Bad labels & A human evaluator mislabels an oak tree as a maple\\
        \bottomrule
      \end{tabular}
    \end{tdmain}
    \begin{tdmargin}
      You can write a program or script to detect omitted values, duplicate
      examples, and out-of-range feature values.
      \par\medskip
      When labels are generated by multiple people, statistically determine
      whether each rater generated equivalent sets of labels. Perhaps one rater
      was a harsher grader, or used different grading criteria.
    \end{tdmargin}
  \end{withmargin}
\end{frame}

\begin{frame}[fragile]{Qualities of good numerical features}
  \begin{withmargin}
    \begin{tdmain}
      \textbf{Clearly named.} Each feature should have a clear, sensible, and
      obvious meaning to any human on the project.
      \vskip0.3em
      \begin{tdcode}
\begin{Verbatim}[fontsize=\scriptsize]
not recommended:  house_age: 851472000
recommended:      house_age_years: 27
\end{Verbatim}
      \end{tdcode}
      \vskip0.4em
      \textbf{Checked or tested before training.} Sometimes bad data, rather
      than bad engineering choices, causes unclear values.
      \vskip0.3em
      \begin{tdcode}
\begin{Verbatim}[fontsize=\scriptsize]
not recommended:  user_age_in_years: 224
recommended:      user_age_in_years: 24
\end{Verbatim}
      \end{tdcode}
    \end{tdmain}
    \begin{tdmargin}
      Although your co-workers will rebel against confusing feature and label
      names, the model won't care, assuming you normalize values properly.
      \par\medskip
      \gterm{Check your data.}
    \end{tdmargin}
  \end{withmargin}
\end{frame}

\begin{frame}[fragile]{Sensible: no magic values}
  \begin{withmargin}
    \begin{tdmain}
      A magic value is a purposeful discontinuity in an otherwise continuous
      feature. A \texttt{watch\_time\_in\_seconds} of $-1$ would force the model
      to figure out what it means to watch a movie \gterm{backwards in time}.
      \vskip0.4em
      \begin{tdcode}
\begin{Verbatim}[fontsize=\scriptsize]
not recommended:
  watch_time_in_seconds: -1

recommended:
  watch_time_in_seconds: 4.82
  is_watch_time_in_seconds_defined = True

  watch_time_in_seconds: 0
  is_watch_time_in_seconds_defined = False
\end{Verbatim}
      \end{tdcode}
    \end{tdmain}
    \begin{tdmargin}
      That is the way to handle a continuous feature with missing values: a
      separate Boolean feature indicating whether a value is supplied.
      \par\medskip
      For a discrete numerical feature like \texttt{product\_category}, signify
      the missing value using a new value in the finite set:
      \par\smallskip
      \texttt{\{0: 'electronics', 1: 'books', 2: 'clothing', 3:
      'missing\_category'\}}
      \par\smallskip
      The model then learns different weights for each value, including missing.
    \end{tdmargin}
  \end{withmargin}
  \keyterms{outliers; feature; feature vector}
\end{frame}

\begin{frame}{Polynomial transforms}
  \begin{withmargin}
    \begin{tdmain}
      Sometimes domain knowledge suggests that one variable is related to the
      square, cube, or other power of another. It's possible to keep the linear
      equation and allow nonlinearity by defining a new term:
      \[ \gterm{x_2} = x_1^2 \qquad\Rightarrow\qquad y = b + w_1x_1 + w_2\gterm{x_2} \]
      \vskip0.4em
      \footnotesize
      This synthetic feature, called a polynomial transform, is treated like any
      other feature. This can still be treated like a linear regression problem,
      and the weights determined through gradient descent as usual, despite
      containing a hidden squared term.
    \end{tdmain}
    \begin{tdmargin}
      Many relationships in the physical world are related to squared terms,
      including acceleration due to gravity, the attenuation of light or sound
      over distance, and elastic potential energy.
      \par\medskip
      If you transform a feature in a way that changes its scale, consider
      normalizing it as well.
      \par\medskip
      The related concept in categorical data is the \textbf{feature cross},
      which more frequently synthesizes two different features.
    \end{tdmargin}
  \end{withmargin}
  \keyterms{categorical data; feature; feature cross; gradient descent; linear
    regression; synthetic feature; weight}
\end{frame}

\begin{frame}{Best practices for numerical data}
  \begin{withmargin}
    \begin{tdmain}
      \footnotesize
      \begin{itemize}
        \item Remember that your model interacts with the data in the feature
          vector, not the data in the dataset.
        \item Normalize most numerical features. If your first strategy doesn't
          succeed, consider a different one.
        \item Binning is sometimes better than normalizing.
        \item Write verification tests to validate what your data should look
          like. The absolute value of latitude should never exceed 90. If your
          data is restricted to Florida, latitudes should fall between 24 and 31.
        \item Visualize your data with scatter plots and histograms. Look for
          anomalies.
        \item Gather statistics not only on the entire dataset but also on
          smaller subsets, because aggregate statistics sometimes obscure
          problems.
        \item Document all your data transformations.
      \end{itemize}
    \end{tdmain}
    \begin{tdmargin}
      \vfill
      {\color{tdfg}\textbf{A model's health is determined by its data.}}
      \par\medskip
      Feed your model healthy data and it will thrive; feed your model junk and
      its predictions will be worthless.
      \par\medskip
      Data is your \gterm{most valuable} resource, so treat it with care.
      \vfill
    \end{tdmargin}
  \end{withmargin}
  \keyterms{binning; bucketing; dataset; feature; feature vector; normalization}
\end{frame}

% ======================================================================
\sectionsubtitle{Encoding a specific set of possible values.}
\section[Categorical data]{Working with Categorical Data}

\begin{frame}{\modulehead{Working with categorical data}{50 minutes}}
  \begin{withmargin}
    \begin{tdmain}
      Categorical data has a specific set of possible values: the different
      species of animals in a national park, the names of streets in a
      particular city, whether or not an email is spam, the colors that house
      exteriors are painted, and binned numbers.
      \vskip0.8em
      True numerical data can be meaningfully \gterm{multiplied}. All else being
      equal, a house of 200 square meters should be roughly twice as valuable as
      an identical house of 100 square meters.
    \end{tdmain}
    \begin{tdmargin}
      Often you should represent features that contain integer values as
      categorical data instead of numerical data.
      \par\medskip
      If you represent postal codes numerically, you're telling the model to
      treat postal code 20004 as twice as large a signal as postal code 10002.
      Representing them as categorical data lets the model weight each
      individual postal code separately.
    \end{tdmargin}
  \end{withmargin}
\end{frame}

\begin{frame}{Vocabulary and index numbers}
  \begin{withmargin}
    \begin{tdmain}
      When a categorical feature has a low number of possible categories, encode
      it as a \gterm{vocabulary}. The model treats each possible categorical
      value as a separate feature and learns different weights for each.
      \vskip0.6em
      \footnotesize
      \begin{tabular}{@{}lrl@{}}
        \toprule
        \textbf{Feature name} & \textbf{Categories} & \textbf{Sample categories}\\
        \midrule
        \texttt{snowed\_today} & 2 & True, False\\
        \texttt{skill\_level} & 3 & Beginner, Practitioner, Expert\\
        \texttt{season} & 4 & Winter, Spring, Summer, Autumn\\
        \texttt{day\_of\_week} & 7 & Monday, Tuesday, Wednesday\\
        \texttt{planet} & 8 & Mercury, Venus, Earth\\
        \bottomrule
      \end{tabular}
    \end{tdmain}
    \begin{tdmargin}
      Machine learning models can only manipulate floating-point numbers, so you
      must convert each string to a unique index number.
      \par\medskip
      But if the data is left as indexed integers and loaded into a model, the
      model would treat the indexed values as continuous floating-point numbers.
      It would consider \texttt{purple} six times more likely than
      \texttt{orange}.
    \end{tdmargin}
  \end{withmargin}
\end{frame}

\begin{frame}{One-hot encoding}
  \begin{withmargin}
    \begin{tdmain}
      Each category is represented by a vector of $N$ elements, where $N$ is the
      number of categories. Exactly one of the elements has the value 1.0; all
      the remaining elements have the value 0.0.
      \vskip0.5em
      \footnotesize
      \begin{tabular}{@{}lrrrrrrrr@{}}
        \toprule
        \textbf{Feature} & Red & Orange & Blue & Yellow & Green & Black & Purple & Brown\\
        \midrule
        "Red"    & 1 & 0 & 0 & 0 & 0 & 0 & 0 & 0\\
        "Orange" & 0 & 1 & 0 & 0 & 0 & 0 & 0 & 0\\
        "Blue"   & 0 & 0 & 1 & 0 & 0 & 0 & 0 & 0\\
        "Yellow" & 0 & 0 & 0 & 1 & 0 & 0 & 0 & 0\\
        "Green"  & 0 & 0 & 0 & 0 & 1 & 0 & 0 & 0\\
        \bottomrule
      \end{tabular}
    \end{tdmain}
    \begin{tdmargin}
      It is the one-hot vector, not the string or the index number, that gets
      passed to the feature vector. The model learns a \gterm{separate weight}
      for each element.
      \par\medskip
      In a true one-hot encoding, only one element has the value 1.0. In a
      variant known as \textbf{multi-hot} encoding, multiple values can be 1.0.
    \end{tdmargin}
  \end{withmargin}
\end{frame}

\begin{frame}[fragile]{Sparse representation}
  \begin{withmargin}
    \begin{tdmain}
      A feature whose values are predominantly zero is a \gterm{sparse feature}.
      Sparse representation means storing the position of the 1.0. The one-hot
      vector for \texttt{"Blue"} is
      \vskip0.3em
      \begin{tdcode}
\begin{Verbatim}[fontsize=\scriptsize]
[0, 0, 1, 0, 0, 0, 0, 0]
\end{Verbatim}
      \end{tdcode}
      \vskip0.3em
      and since the 1 is in position 2, counting from 0, the sparse
      representation is simply \texttt{2}.
    \end{tdmain}
    \begin{tdmargin}
      The sparse representation consumes far less memory than the eight-element
      one-hot vector. Importantly, the model must \gterm{train on the one-hot
      vector}, not the sparse representation.
      \par\medskip
      The sparse representation of a multi-hot encoding stores the positions of
      all the nonzero elements. A car that is both \texttt{"Blue"} and
      \texttt{"Black"} is \texttt{2, 5}.
    \end{tdmargin}
  \end{withmargin}
\end{frame}

\begin{frame}{Outliers and high dimensionality}
  \begin{withmargin}
    \begin{tdmain}
      Rather than giving each rarely used outlier colour such as
      \texttt{"Mauve"} its own category, lump them into a single catch-all
      category called \gterm{out-of-vocabulary} (OOV). The system learns a
      single weight for that outlier bucket.
      \vskip0.6em
      \footnotesize
      \begin{tabular}{@{}lrl@{}}
        \toprule
        \textbf{Feature name} & \textbf{Categories} & \textbf{Sample}\\
        \midrule
        \texttt{words\_in\_english} & $\sim$500{,}000 & "happy", "walking"\\
        \texttt{US\_postal\_codes} & $\sim$42{,}000 & "02114", "90301"\\
        \texttt{last\_names\_in\_Germany} & $\sim$850{,}000 & "Schmidt", "Schneider"\\
        \bottomrule
      \end{tabular}
    \end{tdmain}
    \begin{tdmargin}
      When the number of categories is high, one-hot encoding is usually a bad
      choice. \textbf{Embeddings} are usually much better: the model trains
      faster and infers predictions more quickly, so it has lower latency.
      \par\medskip
      \textbf{Hashing}, the hashing trick, is a less common way to reduce
      dimensions.
    \end{tdmargin}
  \end{withmargin}
\end{frame}

\begin{frame}[fragile]{Hashing}
  \begin{withmargin}
    \begin{tdmain}
      Hashing maps a category to a small integer, the number of the bucket that
      will hold that category.
      \vskip0.4em
      \begin{tdcode}
\begin{Verbatim}[fontsize=\scriptsize]
1. Set the number of bins to N, less than
   the total number of categories.  Say N = 100.
2. Choose a hash function.
3. Pass each category through it: say 89237.
4. Bin index = hash mod N:  89237 % 100 = 37
5. One-hot encode each bin with the new index.
\end{Verbatim}
      \end{tdcode}
    \end{tdmain}
    \begin{tdmargin}
      A machine learning model \gterm{cannot} train directly on raw string
      values like "Red" and "Black". During training, a model can only
      manipulate floating-point numbers.
    \end{tdmargin}
  \end{withmargin}
  \keyterms{categorical data; dimensions; discrete feature; feature cross;
    feature vector; hashing; one-hot encoding; sparse feature; sparse
    representation}
\end{frame}

\begin{frame}{Who labels the categories}
  \begin{withmargin}
    \begin{tdmain}
      Numerical data is often recorded by scientific instruments. Categorical
      data is often categorized by human beings or by ML models. Who decides on
      categories and labels, and how they make those decisions, affects the
      reliability and usefulness of that data.
      \vskip0.7em
      \textbf{Gold labels} are manually labeled by human beings, and are
      considered more desirable due to relatively better data quality. This
      doesn't mean any set of human-labeled data is high quality. Human errors,
      bias, and malice can be introduced at collection or during cleaning.
      \vskip0.5em
      \textbf{Silver labels} are machine-labeled, and can vary widely in
      quality.
    \end{tdmain}
    \begin{tdmargin}
      Any two human beings may label the same example differently. The
      difference between raters' decisions is \gterm{inter-rater agreement}.
      Measure it with Cohen's kappa, intra-class correlation, or Krippendorff's
      alpha.
      \par\medskip
      Check machine labels not only for accuracy and biases but also for
      violations of common sense, reality, and intention: a computer-vision
      model mislabelling a chihuahua as a muffin.
    \end{tdmargin}
  \end{withmargin}
  \slidesources{Hallgren2012}
\end{frame}

\begin{frame}{Feature crosses}
  \begin{withmargin}
    \begin{tdmain}
      Feature crosses are created by taking the Cartesian product of two or more
      categorical or bucketed features. Like polynomial transforms, feature
      crosses allow linear models to handle nonlinearities. They also encode
      \gterm{interactions} between features.
      \vskip0.5em
      \footnotesize
      For a leaf dataset with \texttt{edges} in \{smooth, toothed, lobed\} and
      \texttt{arrangement} in \{opposite, alternate\}, the cross is
      \vskip0.3em
      \texttt{\{Smooth\_Opposite, Smooth\_Alternate, Toothed\_Opposite,\\
      Toothed\_Alternate, Lobed\_Opposite, Lobed\_Alternate\}}
      \vskip0.3em
      where each term is the product of the base feature values. A leaf with a
      lobed edge and alternate arrangement gives \texttt{\{0, 0, 0, 0, 0, 1\}}.
    \end{tdmain}
    \begin{tdmargin}
      Polynomial transforms typically combine numerical data; feature crosses
      combine categorical data.
      \par\medskip
      {\color{tdfg}\textbf{Be careful.}} Crossing two sparse features produces
      an even sparser new feature. A 100-element sparse feature crossed with a
      200-element sparse feature yields a 20{,}000-element sparse feature.
    \end{tdmargin}
  \end{withmargin}
  \keyterms{feature cross; neural network; inter-rater agreement}
\end{frame}

% ======================================================================
\sectionsubtitle{Making good predictions on data the model has never seen.}
\section[Overfitting]{Datasets, Generalization, and Overfitting}

\begin{frame}{\modulehead{Datasets, generalization, and overfitting}{105 minutes}}
  \begin{withmargin}
    \begin{tdmain}
      Two leading questions open this module.
      \vskip0.7em
      \textbf{If you had to prioritize improving one area in your ML project,
      which would have the most impact?} Improving the quality of your dataset.
      \gterm{Data trumps all.} The quality and size of the dataset matters much
      more than which shiny algorithm you use to build your model.
      \vskip0.6em
      \textbf{How much time do you typically spend on data preparation and
      transformation?} More than half the project time. Typically, 80\% of the
      time on a machine learning project is spent constructing datasets and
      transforming data.
    \end{tdmain}
    \begin{tdmargin}
      A better loss function can help a model train faster, but it's still a
      distant second to the quality of the dataset.
      \par\medskip
      This module covers the characteristics of ML datasets, and how to prepare
      your data to ensure high-quality results when training and evaluating your
      model.
    \end{tdmargin}
  \end{withmargin}
\end{frame}

\begin{frame}{Quantity of data}
  \begin{withmargin}
    \begin{tdmain}
      As a rough rule of thumb, your model should train on at least an
      \gterm{order of magnitude} (or two) more examples than trainable
      parameters. However, good models generally train on substantially more
      examples than that.
      \vskip0.8em
      Models trained on large datasets with few features generally outperform
      models trained on small datasets with a lot of features. Google has
      historically had great success training simple models on large datasets.
    \end{tdmain}
    \begin{tdmargin}
      For some relatively simple problems, a few dozen examples might be
      sufficient. For other problems, a trillion examples might be insufficient.
      \par\medskip
      It's possible to get good results from a small dataset if you are adapting
      an existing model already trained on large quantities of data from the
      same schema.
    \end{tdmargin}
  \end{withmargin}
\end{frame}

\begin{frame}{Quality and reliability}
  \begin{withmargin}
    \begin{tdmain}
      \vskip0.3em
      \begin{tdblock}
        A high-quality dataset helps your model accomplish its goal. A low
        quality dataset inhibits your model from accomplishing its goal.
      \end{tdblock}
      \vskip0.6em
      \footnotesize
      \textbf{Reliability} refers to the degree to which you can trust your
      data. In measuring it, determine:
      \vskip0.3em
      \begin{itemize}
        \item How common are label errors? If your data is labeled by humans,
          how often did your raters make mistakes?
        \item Are your features noisy? Be realistic: you can't purge all noise.
          GPS measurements always fluctuate a little, week to week.
        \item Is the data properly filtered for your problem? Should your
          dataset include search queries from bots? For a spam-detection system,
          likely yes. For improving search results for humans, no.
      \end{itemize}
    \end{tdmain}
    \begin{tdmargin}
      {\color{tdfg}\textbf{Common causes of unreliable data.}}
      \par\medskip
      Omitted values. Duplicate examples. Bad feature values. Bad labels. Bad
      sections of data, for example a very reliable feature except for the one
      day when the network kept crashing.
      \par\medskip
      Use \gterm{automation} to flag unreliable data: unit tests that rely on an
      external formal data schema can flag values outside a defined range.
    \end{tdmargin}
  \end{withmargin}
\end{frame}

\begin{frame}{Complete versus incomplete examples}
  \begin{withmargin}
    \begin{tdmain}
      In a perfect world, each example contains a value for each feature.
      Real-world examples are often incomplete. \alertbf{Don't train a model on
      incomplete examples.} Instead:
      \vskip0.5em
      \begin{itemize}
        \item Delete incomplete examples.
        \item Impute missing values, converting the incomplete example to a
          complete one by providing well-reasoned guesses.
      \end{itemize}
      \vskip0.5em
      \footnotesize
      If you can't decide whether to delete or impute, consider building two
      datasets, one formed by deleting and the other by imputing, then determine
      which trains the better model.
    \end{tdmain}
    \begin{tdmargin}
      Imputation is the process of generating well-reasoned data, not random or
      deceptive data. Good imputation can improve your model; bad imputation can
      hurt it.
      \par\medskip
      A good dataset tells the model which values are imputed: add a Boolean
      column such as \texttt{temperature\_is\_imputed}. The model then gradually
      learns to trust imputed examples less.
      \par\medskip
      One common algorithm uses the mean or median. With Z-scores, the imputed
      value is typically 0.
    \end{tdmargin}
  \end{withmargin}
  \keyterms{dataset; embedding vector; example; feature; model; value
    imputation; Z-score normalization}
\end{frame}

\begin{frame}{Direct versus proxy labels}
  \begin{withmargin}
    \begin{tdmain}
      \textbf{Direct labels} are identical to the prediction your model is
      trying to make. A column named \texttt{bicycle owner} is a direct label
      for a model that predicts whether a person owns a bicycle.
      \vskip0.6em
      \textbf{Proxy labels} are similar but not identical. A person subscribing
      to \textit{Bicycle Bizarre} magazine probably, but not definitely, owns a
      bicycle.
      \vskip0.6em
      \footnotesize
      Direct labels are generally better. Proxy labels are always a compromise,
      an imperfect approximation. Models that use proxy labels are only as
      useful as the \gterm{connection} between the proxy label and the
      prediction.
    \end{tdmain}
    \begin{tdmargin}
      Sometimes a direct label exists but can't easily be represented as a
      floating-point number. In that case, use a proxy label.
      \par\medskip
      \texttt{recently bought a bicycle} is a relatively good proxy for bicycle
      ownership. Like all proxy labels it is imperfect: the person buying an
      item isn't always the person using it. People sometimes buy bicycles as a
      gift.
    \end{tdmargin}
  \end{withmargin}
\end{frame}

\begin{frame}{Human-generated data}
  \begin{withmargin}
    \begin{tdmain}
      \textbf{Advantages.} Human raters can perform a wide range of tasks that
      even sophisticated ML models find difficult. The process forces the owner
      of the dataset to develop clear and consistent criteria.
      \vskip0.6em
      \textbf{Disadvantages.} You typically pay human raters, so
      human-generated data can be expensive. To err is human, so multiple raters
      might have to evaluate the same data.
      \vskip0.6em
      \footnotesize
      Always double-check your human raters. Label 1000 examples yourself and
      see how your results match. If discrepancies surface, \gterm{don't assume
      your ratings are the correct ones}, especially if a value judgment is
      involved.
    \end{tdmain}
    \begin{tdmargin}
      {\color{tdfg}\textbf{Questions to settle first.}}
      \par\medskip
      How skilled must your raters be? Must they know a specific language? Do
      you need linguists for dialogue or NLP applications?
      \par\smallskip
      How many labeled examples do you need, and how soon?
      \par\smallskip
      What's your budget?
    \end{tdmargin}
  \end{withmargin}
  \keyterms{label; feature vector}
\end{frame}

\begin{frame}{Class-imbalanced datasets}
  \begin{withmargin}
    \begin{tdmain}
      In a class-balanced dataset, the number of positive and negative classes
      is about equal. In a class-imbalanced dataset, one label is considerably
      more common. \alertbf{In the real world, class-imbalanced datasets are far
      more common.}
      \vskip0.7em
      \footnotesize
      In a dataset of credit card transactions, fraudulent purchases might make
      up less than 0.1\% of examples. In a medical diagnosis dataset, patients
      with a rare virus might be less than 0.01\%. The more common label is the
      \textbf{majority class}; the less common one is the \textbf{minority
      class}.
    \end{tdmain}
    \begin{tdmargin}
      With 200 majority labels and 2 minority labels, a batch size of 20 means
      most batches contain no minority examples at all. A batch size of 100
      averages only one, which is insufficient for proper training.
      \par\medskip
      Accuracy is usually a poor metric for assessing a model trained on a
      class-imbalanced dataset.
    \end{tdmargin}
  \end{withmargin}
\end{frame}

\begin{frame}{Downsample and upweight the majority class}
  \begin{withmargin}
    \begin{tdmain}
      During training, a model should learn two things: what each class looks
      like, and how common each class is. Standard training conflates these two
      goals. This two-step technique separates them.
      \vskip0.6em
      \textbf{Step 1: downsample.} Train on a disproportionately low percentage
      of majority class examples. Downsampling a 99\%/1\% dataset by a factor of
      25 gives a training set that is 80\% majority and 20\% minority.
      \vskip0.5em
      \textbf{Step 2: upweight.} Treat the loss on a majority class example more
      harshly, by the same factor. When the model mistakenly predicts the
      majority class, treat the loss as if it were \gterm{25 errors}.
    \end{tdmain}
    \begin{tdmargin}
      Many students read this section and say some variant of "that just can't
      be right". Be warned that downsampling and upweighting the majority class
      is somewhat counterintuitive.
      \par\medskip
      Downsampling introduces a prediction bias, by showing the model an
      artificial world where the classes are more balanced than in the real
      world. Upweighting corrects it.
      \par\medskip
      Benefits: a better model, which knows both the connection between features
      and labels and the true distribution of the classes, and faster
      convergence.
    \end{tdmargin}
  \end{withmargin}
  \keyterms{batch; class-balanced dataset; class-imbalanced dataset; dataset;
    downsampling; hyperparameter; majority class; minority class; prediction
    bias; upweighting}
\end{frame}

\begin{frame}{Training, validation, and test sets}
  \begin{withmargin}
    \begin{tdmain}
      Splitting into just training and test is a decent idea, but a better
      approach is three subsets. The \gterm{validation set} performs the initial
      testing on the model as it is being trained.
      \vskip0.6em
      A typical split is 70\% training, 15\% validation, 15\% test. Use the
      validation set to evaluate results from the training set. After repeated
      use of the validation set suggests your model is making good predictions,
      use the test set to double-check your model.
      \vskip0.5em
      \footnotesize
      Tweak model means adjusting anything about the model, from changing the
      learning rate, to adding or removing features, to designing a completely
      new model from scratch.
    \end{tdmain}
    \begin{tdmargin}
      {\color{tdfg}\textbf{Why not just two sets.}}
      \par\medskip
      Doing many rounds of tuning against the test set might cause the model to
      implicitly fit the peculiarities of the test set. Like a teacher teaching
      to the test, the model inadvertently fits the test set.
      \par\medskip
      When you transform a feature in your training set, you must make the same
      transformation in the validation set, test set, and real-world dataset.
    \end{tdmargin}
  \end{withmargin}
\end{frame}

\begin{frame}{Test sets wear out}
  \begin{withmargin}
    \begin{tdmain}
      Even with the three-way split, test sets and validation sets still wear
      out with repeated use. The more you use the same data to make decisions
      about hyperparameter settings or other model improvements, the less
      confidence that the model will make good predictions on new data. Collect
      more data to refresh them. \gterm{Starting anew is a great reset.}
      \vskip0.7em
      \footnotesize
      A good test or validation set is large enough to yield statistically
      significant results, representative of the dataset as a whole,
      representative of the real-world data the model will encounter, and has
      \textbf{zero examples duplicated in the training set}.
    \end{tdmain}
    \begin{tdmargin}
      A spam model achieves 99\% precision on both training and test sets. You'd
      expect lower precision on the test set. It turns out many test examples are
      duplicates of training examples: duplicate entries were never scrubbed
      before splitting. You've inadvertently trained on some of your test data.
      \par\medskip
      Dividing examples into train, test, and validation sets is a zero-sum
      game. This is the central trade-off.
    \end{tdmargin}
  \end{withmargin}
  \keyterms{test set; training set; validation set}
\end{frame}

\begin{frame}{Overfitting}
  \begin{withmargin}
    \begin{tdmain}
      Overfitting means creating a model that matches the training set so
      closely that the model fails to make correct predictions on new data. An
      overfit model is analogous to an invention that \alertbf{performs well in
      the lab but is worthless in the real world}.
      \vskip0.8em
      An underfit model doesn't even make good predictions on the training data.
      If an overfit model is like a product that performs well in the lab but
      poorly in the real world, an underfit model is like a product that doesn't
      even do well in the lab.
      \vskip0.5em
      \footnotesize
      Generalization is the opposite of overfitting. Your goal is to create a
      model that generalizes well to new data.
    \end{tdmain}
    \begin{tdmargin}
      In the course's forest example, a complex shape separates all but two of
      the training trees. If we think of the shapes as a model, this is a
      fantastic model. Or is it? On the test set the same model miscategorizes
      many trees.
      \par\medskip
      Overfitting is a common problem in machine learning, not an academic
      hypothetical.
    \end{tdmargin}
  \end{withmargin}
\end{frame}

\begin{frame}{Detecting overfitting}
  \begin{withmargin}
    \begin{tdmain}
      A \textbf{loss curve} plots a model's loss against the number of training
      iterations. A graph that shows two or more loss curves is a
      \gterm{generalization curve}.
      \vskip0.8em
      When the two curves behave similarly at first and then diverge, so that
      loss declines or holds steady for the training set but \textbf{increases}
      for the validation set, that suggests overfitting. In contrast, a
      generalization curve for a well-fit model shows two loss curves that have
      similar shapes.
      \vskip0.6em
      \footnotesize
      Broadly, overfitting is caused by one or both of: the training set doesn't
      adequately represent real life data, or the model is too complex.
    \end{tdmain}
    \begin{tdmargin}
      {\color{tdfg}\textbf{Generalization conditions.}}
      \par\medskip
      Examples must be independently and identically distributed, a fancy way of
      saying your examples can't influence each other.
      \par\smallskip
      The dataset is stationary, meaning it doesn't change significantly over
      time.
      \par\smallskip
      The dataset partitions have the same distribution.
    \end{tdmargin}
  \end{withmargin}
\end{frame}

\begin{frame}{When the conditions break}
  \begin{withmargin}
    \begin{tdmain}
      \textbf{Not stationary.} A streaming service predicting the popularity of
      new shows for the next three years, trained on ten years of data, will
      struggle: viewer tastes change in ways that past behavior can't predict.
      \vskip0.6em
      \textbf{Not independent.} A model recommending the ideal date to buy a
      train ticket exhibits low loss on validation and test sets but sometimes
      makes terrible real-world predictions, because of a
      \alertbf{feedback loop}. Riders who follow the recommendation buy at 8:30
      on July 8; at 9:00 the company raises prices because fewer seats remain.
      Riders using the model's recommendation have altered prices.
    \end{tdmain}
    \begin{tdmargin}
      To make partitions statistically similar, shuffle the examples in the
      dataset extensively before partitioning them.
      \par\medskip
      Sorting examples from earliest to most recent makes the partitions less
      similar if the examples are not stationary.
      \par\medskip
      Given enough examples, the law of averages does \textbf{not} naturally
      ensure similar distributions.
    \end{tdmargin}
  \end{withmargin}
  \keyterms{feedback loop; generalization; generalization curve; independently
    and identically distributed; loss curve; overfitting; stationarity; training
    set; underfitting}
\end{frame}

\begin{frame}{Occam's razor}
  \begin{withmargin}
    \begin{tdmain}
      Replace the complex forest model with a ridiculously simple one, a
      straight line, and it \gterm{generalizes better} on new data.
      \vskip0.8em
      Simplicity has been beating complexity for a long time. The preference for
      simplicity dates back to ancient Greece. Centuries later, a
      fourteenth-century friar named William of Occam formalized it in a
      philosophy known as Occam's razor. This philosophy remains an essential
      underlying principle of many sciences, including machine learning.
      \vskip0.5em
      \footnotesize
      Complex models typically outperform simple models on the training set.
      However, simple models typically outperform complex models on the test
      set, which is more important.
    \end{tdmain}
    \begin{tdmargin}
      {\color{tdfg}\textbf{How many features to start with?}}
      \par\medskip
      Pick 1 to 3 that seem to have strong predictive power. Every new feature
      adds a dimension. When the dimensionality increases, the volume of the
      space increases so fast that the available training data become sparse.
      This is the \textbf{curse of dimensionality}.
      \par\medskip
      The two most famous physics formulas of all time, $F=ma$ and $E=mc^2$,
      each involve only three variables.
    \end{tdmargin}
  \end{withmargin}
\end{frame}

\begin{frame}{Loss plus complexity}
  \begin{withmargin}
    \begin{tdmain}
      Machine learning models must simultaneously meet two conflicting goals:
      fit data well, and fit data as simply as possible. One approach is to
      penalize complex models, forcing the model to become simpler during
      training. That is \gterm{regularization}.
      \[ \text{minimize(loss)} \quad\longrightarrow\quad \text{minimize(loss} + \gterm{\text{complexity}}) \]
      \vskip0.3em
      \footnotesize
      Loss and complexity are typically inversely related. As complexity
      increases, loss decreases. Find a reasonable middle ground where the model
      makes good predictions on both the training data and real-world data.
    \end{tdmain}
    \begin{tdmargin}
      {\color{tdfg}\textbf{A regularization analogy.}}
      \par\medskip
      Suppose every student in a lecture hall had a buzzer that annoyed the
      professor, pressed whenever the lecture became too complicated. The
      professor would complain, "when I simplify, I'm not being precise enough".
      The students would counter, "the only goal is to explain it simply enough
      that I understand it". Gradually the buzzers would train the professor to
      give an appropriately simple lecture.
    \end{tdmargin}
  \end{withmargin}
\end{frame}

\begin{frame}{L2 regularization}
  \begin{withmargin}
    \begin{tdmain}
      \[ L_2 \text{ regularization} = w_1^2 + w_2^2 + \ldots + w_n^2 \]
      \vskip0.4em
      \footnotesize
      \begin{tabular}{@{}lrr@{}}
        \toprule
         & \textbf{Value} & \textbf{Squared value}\\
        \midrule
        $w_1$ & 0.2  & 0.04\\
        $w_2$ & -0.5 & 0.25\\
        $w_3$ & 5.0  & \gterm{25.0}\\
        $w_4$ & -1.2 & 1.44\\
        $w_5$ & 0.3  & 0.09\\
        $w_6$ & -0.1 & 0.01\\
        \midrule
         &  & 26.83 total\\
        \bottomrule
      \end{tabular}
    \end{tdmain}
    \begin{tdmargin}
      Weights close to zero don't affect L2 regularization much, but large
      weights have a huge impact. A single weight, $w_3$, contributes about 93\%
      of the total complexity. The other five contribute only about 7\%.
      \par\medskip
      L2 regularization encourages weights toward 0, but \textbf{never} pushes
      weights all the way to zero. Consequently, all features still contribute
      something to the model.
    \end{tdmargin}
  \end{withmargin}
\end{frame}

\begin{frame}{Regularization rate (lambda)}
  \begin{withmargin}
    \begin{tdmain}
      \[ \text{minimize(loss} + \gterm{\lambda}\ \text{complexity)} \]
      \vskip0.5em
      \textbf{A high regularization rate} strengthens the influence of
      regularization, reducing the chances of overfitting. It tends to produce a
      histogram of model weights that is a normal distribution with a mean
      weight of 0.
      \vskip0.5em
      \textbf{A low regularization rate} lowers the influence of regularization,
      increasing the chances of overfitting, and tends to produce a flat
      distribution of weights.
    \end{tdmain}
    \begin{tdmargin}
      Setting the regularization rate to zero removes regularization completely.
      Training then focuses exclusively on minimizing loss, which poses the
      highest possible overfitting risk.
      \par\medskip
      The ideal regularization rate produces a model that generalizes well to
      new, previously unseen data. That ideal value is data-dependent, so you
      must do some tuning.
    \end{tdmargin}
  \end{withmargin}
\end{frame}

\begin{frame}{Early stopping, and the equilibrium}
  \begin{withmargin}
    \begin{tdmain}
      \textbf{Early stopping} is a regularization method that doesn't involve a
      calculation of complexity. It simply means ending training before the
      model fully converges, for example when the loss curve for the validation
      set starts to increase. It usually increases training loss but can
      decrease test loss. It is a quick, but rarely optimal, form of
      regularization.
      \vskip0.8em
      Learning rate and regularization rate tend to move weights in
      \gterm{opposite directions}. A high learning rate often pulls weights away
      from zero; a high regularization rate pushes weights towards zero.
    \end{tdmain}
    \begin{tdmargin}
      If the regularization rate is high with respect to the learning rate, the
      weak weights tend to produce a model that makes poor predictions.
      \par\medskip
      If the learning rate is high with respect to the regularization rate, the
      strong weights tend to produce an overfit model.
      \par\medskip
      Worst of all, once you find that elusive balance, you may ultimately have
      to change the learning rate. And then you'll again have to find the ideal
      regularization rate.
    \end{tdmargin}
  \end{withmargin}
  \keyterms{early stopping; L1 regularization; L2 regularization; learning rate;
    regularization; regularization rate}
\end{frame}

\begin{frame}{Reading four broken loss curves}
  \begin{withmargin}
    \begin{tdmain}
      \footnotesize
      \begin{tabular}{@{}p{0.34\linewidth}p{0.60\linewidth}@{}}
        \toprule
        \textbf{Symptom} & \textbf{What to do}\\
        \midrule
        Loss oscillates erratically, never flattening
          & Check your data against a data schema and remove bad examples.
            Reduce the learning rate. Reduce the training set to a tiny number
            of trustworthy examples, then gradually add more.\\[0.4em]
        Loss decreases, then a sharp jump up
          & The input data contains one or more NaNs, for example from a
            division by zero. Or a batch contains a burst of outliers, due to
            improper shuffling.\\[0.4em]
        Training loss converges, validation loss rises
          & The model is overfitting. Make the model simpler, increase the
            regularization rate, or ensure the training and test sets are
            statistically equivalent.\\[0.4em]
        Loss converges, then repeats a rectangular-wave pattern
          & The training set is not shuffled well. 100 images of dogs followed
            by 100 images of cats will make loss oscillate.\\
        \bottomrule
      \end{tabular}
    \end{tdmain}
    \begin{tdmargin}
      \gterm{Avoid increasing the learning rate} when a model's learning curve
      indicates a problem.
      \par\medskip
      Increasing the number of examples in the training set is a tempting idea
      for the oscillating curve, but is extremely unlikely to fix the problem.
      \par\medskip
      A very high regularization rate could prevent a model from converging, but
      won't cause a sharp jump.
    \end{tdmargin}
  \end{withmargin}
  \keyterms{batch; example; feature; learning rate; loss curve; outliers;
    overfitting; regularization; regularization rate; test set; training set;
    validation set}
\end{frame}

% ======================================================================
\sectionsubtitle{Learning the feature crosses for you.}
\section[Neural networks]{Neural Networks}

\begin{frame}{\modulehead{Neural networks}{75 minutes}}
  \begin{withmargin}
    \begin{tdmain}
      Nonlinear means that you can't accurately predict a label with a model of
      the form $b + w_1x_1 + w_2x_2$. In other words, the decision surface is
      not a line.
      \vskip0.8em
      Adding a feature cross $x_3 = x_1x_2$ lets a linear model learn a
      hyperbolic shape. But determining the correct feature crosses takes effort
      and experimentation. What if you didn't have to do all that yourself?
      \vskip0.6em
      Neural networks are a family of model architectures designed to find
      nonlinear patterns in data. During training, the model
      \gterm{automatically learns} the optimal feature crosses to perform on the
      input data to minimize loss.
    \end{tdmain}
    \begin{tdmargin}
      {\color{tdfg}\textbf{Objectives.}}
      \par\medskip
      Explain the motivation for building neural networks. Define nodes, hidden
      layers, and activation functions. Step through the inference process.
      Build a high-level intuition of backpropagation. Explain one-vs.-all and
      one-vs.-one multi-class classification.
    \end{tdmargin}
  \end{withmargin}
\end{frame}

\begin{frame}{Nodes and hidden layers}
  \begin{withmargin}
    \begin{tdmain}
      Start with a linear model $y' = b + w_1x_1 + w_2x_2 + w_3x_3$: three input
      nodes and one output node. Adjusting the weights and bias does not change
      the overall mathematical relationship. It is still a linear model.
      \vskip0.7em
      Additional layers between the input layer and the output layer are called
      \gterm{hidden layers}, and the nodes in these layers are called
      \textbf{neurons}. The value of each neuron is calculated the same way as
      the output of a linear model: the sum of the product of each of its inputs
      and a unique weight parameter, plus the bias.
    \end{tdmain}
    \begin{tdmargin}
      {\color{tdfg}\textbf{Counting parameters.}}
      \par\medskip
      One hidden layer of 4 neurons over 3 inputs: 4 parameters per neuron, 3
      weights and a bias, so 16. Then 5 more for the output: 4 weights and a
      bias. In total \textbf{21} parameters.
      \par\medskip
      Can this model learn nonlinearities? \textbf{No.} All the calculations are
      linear, and linear calculations performed on the output of linear
      calculations are also linear.
    \end{tdmargin}
  \end{withmargin}
  \keyterms{bias; hidden layer; input layer; linear model; model; neural
    network; neuron; nonlinear; parameter; weight}
\end{frame}

\begin{frame}{Activation functions}
  \begin{withmargin}
    \begin{tdmain}
      We need some way to insert nonlinear mathematical operations into a model.
      An \gterm{activation function} is a nonlinear transform of a neuron's
      output value before the value is passed as input to the next layer.
      \vskip0.8em
      Now that we've added an activation function, adding layers has more
      impact. Stacking nonlinearities on nonlinearities lets us model very
      complicated relationships between the inputs and the predicted outputs.
      Each layer is effectively learning a more complex, higher-level function
      over the raw inputs.
    \end{tdmain}
    \begin{tdmargin}
      We have already applied a nonlinear operation to a linear model's output:
      in logistic regression, the sigmoid function.
      \par\medskip
      Applying the same principle here, each node's output becomes a sigmoid
      transform of the linear combination of the nodes in the previous layer,
      squished between 0 and 1.
    \end{tdmargin}
  \end{withmargin}
\end{frame}

\begin{frame}{The three common activation functions}
  \begin{withmargin}
    \begin{tdmain}
      \begin{align*}
        \text{sigmoid} \quad F(x) &= \frac{1}{1 + e^{-x}} &&\text{range } 0 \text{ to } 1\\[0.3em]
        \text{tanh} \quad F(x) &= \tanh(x) &&\text{range } -1 \text{ to } 1\\[0.3em]
        \gterm{\text{ReLU}} \quad F(x) &= \max(0, x) &&\text{0 below zero, } x \text{ above}
      \end{align*}
      \vskip0.4em
      \footnotesize
      ReLU often works a little better than a smooth function like sigmoid or
      tanh, because it is less susceptible to the vanishing gradient problem
      during training. ReLU is also significantly easier to compute.
    \end{tdmain}
    \begin{tdmargin}
      The term sigmoid is often used more generally to refer to any S-shaped
      function. A more technically precise term for
      $F(x)=\frac{1}{1+e^{-x}}$ is \textbf{logistic function}.
      \par\medskip
      In practice, any mathematical function can serve as an activation
      function. With $\sigma$ the activation function, the value of a node is
      $\sigma(\boldsymbol{w} \cdot \boldsymbol{x} + b)$.
      \par\medskip
      Keras supports many activation functions out of the box. Still, start with
      ReLU.
    \end{tdmargin}
  \end{withmargin}
\end{frame}

\begin{frame}{What a neural network is}
  \begin{withmargin}
    \begin{tdmain}
      The model now has all the standard components of what people usually mean
      when they refer to a neural network.
      \vskip0.6em
      \begin{itemize}
        \item A set of nodes, analogous to neurons, organized in layers.
        \item A set of learned weights and biases representing the connections
          between each layer and the layer beneath it.
        \item An activation function that transforms the output of each node in
          a layer. Different layers may have different activation functions.
      \end{itemize}
    \end{tdmain}
    \begin{tdmargin}
      {\color{tdfg}\textbf{A caveat.}}
      \par\medskip
      Neural networks aren't necessarily always better than feature crosses, but
      they do offer a \gterm{flexible alternative} that works well in many
      cases.
    \end{tdmargin}
  \end{withmargin}
  \keyterms{activation function; sigmoid function; vanishing gradient problem}
\end{frame}

\begin{frame}{Backpropagation}
  \begin{withmargin}
    \begin{tdmain}
      Backpropagation is the most common training algorithm for neural networks.
      It makes gradient descent feasible for multi-layer neural networks. Many
      ML code libraries handle backpropagation automatically, so you don't need
      to perform any of the underlying calculations yourself.
      \vskip0.8em
      \textbf{Vanishing gradients.} The gradients for the lower layers, those
      closer to the input, can become very small. In deep networks, computing
      these gradients can involve taking the product of many small terms. When
      the values approach 0, the gradients are said to \gterm{vanish}, and those
      layers train very slowly or not at all. ReLU can help prevent this.
    \end{tdmain}
    \begin{tdmargin}
      {\color{tdfg}\textbf{Exploding gradients.}}
      \par\medskip
      If the weights are very large, the gradients for the lower layers involve
      products of many large terms and get too large to converge. Batch
      normalization can help, as can lowering the learning rate.
      \par\medskip
      The backpropagation algorithm uses the calculus concept of a gradient.
      Understanding and debugging these issues usually requires some background
      in calculus.
    \end{tdmargin}
  \end{withmargin}
\end{frame}

\begin{frame}{Dead ReLU units, and dropout}
  \begin{withmargin}
    \begin{tdmain}
      Once the weighted sum for a ReLU unit falls below 0, the unit can get
      stuck. It outputs 0, contributing nothing to the network's output, and
      gradients can no longer flow through it. With a source of gradients cut
      off, the input may never change enough to bring the weighted sum back
      above 0. \alertbf{Lowering the learning rate can help keep ReLU units from
      dying.}
      \vskip0.7em
      \textbf{Dropout regularization} works by randomly dropping out unit
      activations in a network for a single gradient step. The more you drop
      out, the stronger the regularization.
    \end{tdmain}
    \begin{tdmargin}
      0.0 = no dropout regularization.
      \par\smallskip
      1.0 = drop out all nodes. The model learns nothing.
      \par\smallskip
      Values between 0.0 and 1.0 = more useful.
      \par\medskip
      There are also many variants of ReLU designed to address dead units, such
      as LeakyReLU.
    \end{tdmargin}
  \end{withmargin}
  \keyterms{backpropagation; dropout regularization; exploding gradient problem;
    vanishing gradient problem}
\end{frame}

\begin{frame}{One versus all}
  \begin{withmargin}
    \begin{tdmain}
      One-vs.-all provides a way to use binary classification for a series of
      yes or no predictions across multiple possible labels. Given $N$ possible
      solutions, it consists of $N$ separate binary classifiers, one for each
      possible outcome.
      \vskip0.8em
      Given a picture of a piece of fruit, four different recognizers might be
      trained: is this image an apple? an orange? a banana? a grape? This
      approach is reasonable when the total number of classes is small, but
      becomes \gterm{increasingly inefficient} as the number of classes rises.
    \end{tdmain}
    \begin{tdmargin}
      A significantly more efficient one-vs.-all model uses a deep neural
      network in which each output node represents a different class, with a
      sigmoid activation function applied to the output layer.
      \par\medskip
      Some real-world multi-class problems entail choosing from millions of
      separate classes.
    \end{tdmargin}
  \end{withmargin}
\end{frame}

\begin{frame}{One versus one, or softmax}
  \begin{withmargin}
    \begin{tdmain}
      With sigmoid on the output layer, the probabilities \gterm{don't sum to
      1.0}: for the fruit example they sum to 1.43. In a one-vs.-all approach,
      the probability of each binary set of outcomes is determined independently
      of all the others.
      \[ p(y = j|\textbf{x}) = \frac{e^{(\textbf{w}_j^{T}\textbf{x} + b_j)}}
                                    {\sum_{k \in K} e^{(\textbf{w}_k^{T}\textbf{x} + b_k)}} \]
      \vskip0.3em
      \footnotesize
      Softmax assigns decimal probabilities to each class such that all
      probabilities add up to 1.0. This additional constraint helps training
      converge more quickly. The formula extends logistic regression to multiple
      classes.
    \end{tdmain}
    \begin{tdmargin}
      To perform softmax, the hidden layer directly preceding the output layer,
      called the \textbf{softmax layer}, must have the same number of nodes as
      the output layer.
      \par\medskip
      In the course's example, sigmoid predicts a 0.84 chance the image is a
      pear; softmax predicts 0.63, with all four fruits summing to 1.0.
    \end{tdmargin}
  \end{withmargin}
\end{frame}

\begin{frame}{Softmax options, and its one assumption}
  \begin{withmargin}
    \begin{tdmain}
      \textbf{Full softmax} calculates a probability for every possible class.
      It is fairly cheap when the number of classes is small but becomes
      prohibitively expensive when the number climbs.
      \vskip0.6em
      \textbf{Candidate sampling} calculates a probability for all the positive
      labels but only for a random sample of negative labels. If we are
      determining whether an input image is a beagle or a bloodhound, we don't
      have to provide probabilities for every non-doggy example.
      \vskip0.6em
      \footnotesize
      Softmax assumes each example is a member of \gterm{exactly one} class. If
      an input image might contain a bowl of both apples and oranges, you may
      not use softmax; you must rely on multiple logistic regressions.
    \end{tdmain}
    \begin{tdmargin}
      Multi-class examples: is this dog a beagle, a basset hound, or a
      bloodhound? Is that plane a Boeing 747, Airbus 320, Boeing 777, or Embraer
      190?
    \end{tdmargin}
  \end{withmargin}
  \keyterms{binary classification; multi-class classification; one-vs.-all
    classification; softmax}
\end{frame}

% ======================================================================
\sectionsubtitle{Lower-dimensional representations of sparse data.}
\section[Embeddings]{Embeddings}

\begin{frame}{\modulehead{Embeddings}{45 minutes}}
  \begin{withmargin}
    \begin{tdmain}
      Imagine developing a food-recommendation application, where users input
      their favorite meals and the app suggests similar meals they might like.
      You curate a dataset of 5{,}000 popular meal items, including borscht, hot
      dog, salad, pizza, and shawarma, and create a \texttt{meal} feature that
      one-hot encodes each of them.
      \vskip0.8em
      \gterm{Encoding} refers to the process of choosing an initial numerical
      representation of data to train the model on. Each one-hot vector has a
      length of 5{,}000, one entry per menu item.
    \end{tdmain}
    \begin{tdmargin}
      {\color{tdfg}\textbf{Objectives.}}
      \par\medskip
      Visualize vector representations of word embeddings, such as word2vec.
      Distinguish encoding from embedding. Describe contextual embedding.
    \end{tdmargin}
  \end{withmargin}
\end{frame}

\begin{frame}{Pitfalls of sparse data representations}
  \begin{withmargin}
    \begin{tdmain}
      \footnotesize
      \begin{itemize}
        \item \textbf{Number of weights.} With $M$ entries in your one-hot
          encoding and $N$ nodes in the first layer, the model has to train
          $M \times N$ weights for that layer.
        \item \textbf{Number of datapoints.} The more weights, the more data you
          need to train effectively.
        \item \textbf{Amount of computation.} The more weights, the more
          computation required to train and use the model. It's easy to exceed
          the capabilities of your hardware.
        \item \textbf{Amount of memory.} The more weights, the more memory
          needed on the accelerators that train and serve it.
        \item \textbf{On-device ML.} To run on local devices, you need to make
          your model smaller and decrease the number of weights.
      \end{itemize}
    \end{tdmain}
    \begin{tdmargin}
      \vfill
      Embeddings are \gterm{lower-dimensional representations} of sparse data
      that address all of these issues.
      \vfill
    \end{tdmargin}
  \end{withmargin}
  \keyterms{one-hot encoding; neural network; weight}
\end{frame}

\begin{frame}{Embedding space}
  \begin{withmargin}
    \begin{tdmain}
      An embedding is a vector representation of data in embedding space.
      Generally, a model finds potential embeddings by projecting the
      high-dimensional space of initial data vectors into a
      \gterm{lower-dimensional} space.
      \vskip0.8em
      To give an idea, place hot dog, pizza, salad, shawarma, and borscht on a
      single imaginary dimension of "sandwichness". Where would an apple strudel
      fall? Between hot dog and shawarma, but it also has an additional
      dimension of "dessertness". Borscht is much more liquid than the others,
      which suggests a third dimension, "liquidness".
    \end{tdmain}
    \begin{tdmargin}
      An embedding represents each item in $n$-dimensional space with $n$
      floating-point numbers, typically in the range $-1$ to 1 or 0 to 1.
      \par\medskip
      In 2D, apple strudel could be $(0.5, 0.3)$ and hot dog $(0.2, -0.5)$.
      \par\medskip
      The distance between any two items can be calculated mathematically and
      interpreted as a measure of relative \textbf{similarity}.
    \end{tdmargin}
  \end{withmargin}
\end{frame}

\begin{frame}{Real-world embedding spaces}
  \begin{withmargin}
    \begin{tdmain}
      In the real world, embedding spaces are $d$-dimensional, where $d$ is much
      higher than 3 though lower than the dimensionality of the data. For word
      embeddings, $d$ is often 256, 512, or 1024. The individual dimensions are
      rarely as understandable as "dessertness" or "liquidness". Sometimes what
      they mean can be inferred, but \alertbf{not always}.
      \vskip0.8em
      Embeddings will usually be specific to the task. "Cereal" and "breakfast
      sausage" would probably be close together in the embedding space of a
      time-of-day model but far apart in the embedding space of a vegetarian
      versus non-vegetarian model.
    \end{tdmain}
    \begin{tdmargin}
      The ML practitioner usually sets the specific task and the number of
      embedding dimensions. The model then tries to arrange the training
      examples to be close in an embedding space with that number of dimensions,
      or tunes for the number if $d$ is not fixed.
    \end{tdmargin}
  \end{withmargin}
  \slidesources{Chollet2017}
\end{frame}

\begin{frame}{Static embeddings}
  \begin{withmargin}
    \begin{tdmain}
      One task has general applicability: \gterm{predicting the context of a
      word}. Models trained to do so assume that words appearing in similar
      contexts are semantically related. Training data that includes "They rode
      a burro down into the Grand Canyon" and "They rode a horse down into the
      canyon" suggests that "horse" appears in similar contexts to "burro".
      \vskip0.8em
      \texttt{word2vec} trains on a corpus of documents to obtain a single
      global embedding per word. When each word or data point has a single
      embedding vector, this is called a \textbf{static embedding}.
    \end{tdmain}
    \begin{tdmargin}
      Research suggests that static embeddings, once trained, encode some degree
      of semantic information, particularly in relationships between words.
      \par\medskip
      The specific embedding vectors generated depend on the corpus used for
      training.
      \par\medskip
      word2vec refers both to an algorithm for obtaining static word embeddings
      and to a set of pretrained word vectors.
    \end{tdmargin}
  \end{withmargin}
  \slidesources{Mikolov2013}
  \keyterms{embedding vector; embedding space}
\end{frame}

\begin{frame}{Obtaining embeddings}
  \begin{withmargin}
    \begin{tdmain}
      \textbf{Dimensionality reduction.} Principal component analysis has been
      used to create word embeddings. Given a set of instances like bag of words
      vectors, PCA tries to find highly correlated dimensions that can be
      collapsed into a single dimension.
      \vskip0.7em
      \textbf{Training an embedding as part of a neural network.} Create a
      hidden layer of size $d$ designated as the \gterm{embedding layer}, where
      $d$ is both the number of nodes in the hidden layer and the number of
      dimensions in the embedding space. As in any deep network, the parameters
      are optimized during training to minimize loss.
    \end{tdmain}
    \begin{tdmargin}
      For the food recommender: compile users' top five favorite foods, set four
      of them as feature data, and randomly set aside the fifth as the positive
      label the model aims to predict, optimizing with a softmax loss.
      \par\medskip
      A one-hot encoding of \texttt{hot dog} might be translated into the
      three-dimensional embedding vector \texttt{[2.98, -0.75, 0]}.
    \end{tdmargin}
  \end{withmargin}
\end{frame}

\begin{frame}{Contextual embeddings}
  \begin{withmargin}
    \begin{tdmain}
      Words can mean different things in different contexts. "Yeah" means one
      thing on its own, but the opposite in the phrase "Yeah, right." "Post" can
      mean mail, to put in the mail, earring backing, marker at the end of a
      horse race, postproduction, pillar, to put up a notice, to station a
      guard, or after.
      \vskip0.8em
      With static embeddings each word is a \gterm{single point} in vector
      space, even though it may have a variety of meanings. Contextual
      embeddings allow a word to be represented by multiple embeddings that
      incorporate information about the surrounding words as well as the word
      itself.
    \end{tdmain}
    \begin{tdmargin}
      \textbf{ELMo} aggregates the static embedding with embeddings from other
      layers, which encode front-to-back and back-to-front readings of the
      sentence.
      \par\smallskip
      \textbf{BERT} masks part of the input sequence.
      \par\smallskip
      \textbf{Transformers} use a self-attention layer to weight the relevance
      of the other words in a sequence to each individual word, and add a
      positional embedding.
    \end{tdmargin}
  \end{withmargin}
  \keyterms{bag of words; BERT; embedding layer; embedding vector; neural
    network; positional encoding; softmax}
\end{frame}

% ======================================================================
\sectionsubtitle{From tokens to Transformers.}
\section[LLMs]{Intro to Large Language Models}

\begin{frame}[fragile]{\modulehead{Intro to large language models}{45 minutes}}
  \begin{withmargin}
    \begin{tdmain}
      A language model estimates the probability of a token or sequence of
      tokens occurring within a longer sequence of tokens. A token could be a
      word, a subword, or even a single character.
      \vskip0.4em
      \begin{tdcode}
\begin{Verbatim}[fontsize=\scriptsize]
When I hear rain on my roof, I _____ in my kitchen.
\end{Verbatim}
      \end{tdcode}
      \vskip0.4em
      \footnotesize
      \begin{tabular}{@{}rl@{}}
        \toprule
        \textbf{Probability} & \textbf{Token(s)}\\
        \midrule
        9.4\% & cook soup\\
        5.2\% & warm up a kettle\\
        3.6\% & cower\\
        2.5\% & nap\\
        2.2\% & relax\\
        \bottomrule
      \end{tabular}
    \end{tdmain}
    \begin{tdmargin}
      An application can use the probability table to make predictions: the
      highest probability, or a random selection from tokens above a threshold.
      \par\medskip
      By modeling the statistical patterns of tokens, modern language models
      develop extremely powerful internal representations of language and can
      \gterm{generate plausible language}.
    \end{tdmargin}
  \end{withmargin}
\end{frame}

\begin{frame}{Tokens}
  \begin{withmargin}
    \begin{tdmain}
      Most modern language models tokenize by \gterm{subwords}, chunks of text
      containing semantic meaning. The chunks vary in length from single
      characters like punctuation or the possessive s to whole words. Prefixes
      and suffixes might be represented as separate subwords.
      \vskip0.6em
      \footnotesize
      \begin{itemize}
        \item \texttt{unwatched} becomes \texttt{un}, \texttt{watch},
          \texttt{ed}.
        \item \texttt{cats} becomes \texttt{cat}, \texttt{s}.
        \item \texttt{antidisestablishmentarianism} becomes \texttt{anti},
          \texttt{dis}, \texttt{establish}, \texttt{ment}, \texttt{arian},
          \texttt{ism}.
      \end{itemize}
    \end{tdmain}
    \begin{tdmargin}
      Tokenization is language specific, so the number of characters per token
      differs across languages. For English, one token corresponds to about 4
      characters or about three quarters of a word, so 400 tokens is about 300
      English words.
      \par\medskip
      Tokens are the atomic unit of language modeling, and are now also being
      successfully applied to computer vision and audio generation.
    \end{tdmargin}
  \end{withmargin}
\end{frame}

\begin{frame}[fragile]{N-gram language models}
  \begin{withmargin}
    \begin{tdmain}
      N-grams are ordered sequences of words used to build language models,
      where N is the number of words in the sequence.
      \vskip0.4em
      \begin{tdcode}
\begin{Verbatim}[fontsize=\scriptsize]
"you are very nice"

2-grams:  you are / are very / very nice
3-grams:  you are very / are very nice
\end{Verbatim}
      \end{tdcode}
      \vskip0.4em
      \footnotesize
      Given two words as input, a 3-gram model predicts the likelihood of the
      third. Given \texttt{orange is}, hundreds of 3-grams could follow. Two of
      them are \texttt{orange is ripe}, about the fruit, and \texttt{orange is
      cheerful}, about the colour.
    \end{tdmain}
    \begin{tdmargin}
      In language models, \gterm{context} is helpful information before or after
      the target token. Context can help a language model determine whether
      orange refers to a citrus fruit or a colour.
      \par\medskip
      Humans retain relatively long contexts. While watching Act 3 of a play,
      you retain knowledge of characters introduced in Act 1.
    \end{tdmargin}
  \end{withmargin}
\end{frame}

\begin{frame}{Why N-grams run out of context}
  \begin{withmargin}
    \begin{tdmain}
      The only context a 3-gram provides is the first two words. Due to lack of
      context, language models based on 3-grams make a lot of mistakes.
      \vskip0.8em
      Longer N-grams provide more context. However, as N grows, the relative
      occurrence of each instance \gterm{decreases}. When N becomes very large,
      the language model typically has only a single instance of each occurrence
      of N tokens, which isn't very helpful in predicting the target token.
      \vskip0.6em
      \footnotesize
      Is a 6-gram model better than a 5-gram model for English? It depends on
      the size and diversity of the training set. If the training set spans
      millions of diverse documents, the 6-gram model will probably win.
      Otherwise most of the 6-grams will be rare.
    \end{tdmain}
    \begin{tdmargin}
      \textbf{Recurrent neural networks} train on a sequence of tokens and
      provide more context than N-grams. An RNN can gradually learn, and learn
      to ignore, selected context from each word in a sentence, kind of like you
      would when listening to someone speak.
      \par\medskip
      RNNs evaluate information token by token, and training them for long
      contexts is constrained by the vanishing gradient problem.
    \end{tdmargin}
  \end{withmargin}
  \keyterms{bigram; language model; N-gram; neural network; recurrent neural
    network; token; trigram; vanishing gradient problem}
\end{frame}

\begin{frame}{Transformers}
  \begin{withmargin}
    \begin{tdmain}
      LLMs make much better predictions than N-gram language models or recurrent
      neural networks because they contain far more parameters and gather far
      more context. The most successful and widely used architecture for
      building LLMs is the \gterm{Transformer}.
      \vskip0.8em
      A full Transformer consists of an \textbf{encoder}, which converts input
      text into an intermediate representation, and a \textbf{decoder}, which
      converts that intermediate representation into useful text. Both are
      enormous neural nets.
    \end{tdmain}
    \begin{tdmargin}
      \textbf{Encoder-only} architectures map input text into an intermediate
      representation, often an embedding layer. Use cases: predicting any token
      in the input sequence, or creating a sophisticated embedding to serve as
      input for another system.
      \par\medskip
      \textbf{Decoder-only} architectures generate new tokens from the text
      already generated, and typically excel at generating sequences.
    \end{tdmargin}
  \end{withmargin}
  \slidesources{Vaswani2017}
\end{frame}

\begin{frame}[fragile]{Self-attention}
  \begin{withmargin}
    \begin{tdmain}
      On behalf of each token of input, self-attention asks:
      \gterm{how much does each other token of input affect the interpretation
      of this token?}
      \vskip0.4em
      \begin{tdcode}
\begin{Verbatim}[fontsize=\scriptsize]
The animal didn't cross the street
because it was too tired.
\end{Verbatim}
      \end{tdcode}
      \vskip0.4em
      \footnotesize
      Each of the eleven words is paying attention to the other ten. The pronoun
      \texttt{it} typically refers to a recent noun, but which one, the animal
      or the street? Self-attention rates \texttt{animal} as more important.
      Change the final word to \texttt{wide}, and self-attention would hopefully
      rate \texttt{street} as more relevant instead.
    \end{tdmain}
    \begin{tdmargin}
      The "self" refers to the input sequence. Some attention mechanisms weigh
      relations of input tokens to tokens in an output sequence; self-attention
      only weighs relations between tokens in the input sequence.
      \par\medskip
      \textbf{Bidirectional} self-attention gathers context from words on either
      side; encoders use it. \textbf{Unidirectional} gathers from one side only;
      decoders use it, because generating token by token requires it.
    \end{tdmargin}
  \end{withmargin}
\end{frame}

\begin{frame}[fragile]{Multi-head, multi-layer}
  \begin{withmargin}
    \begin{tdmain}
      Each self-attention layer is typically comprised of multiple
      self-attention \textbf{heads}. Since the parameters of each head are
      initialized to random values, different heads can learn different
      relationships. One head might determine which noun a pronoun refers to;
      another might learn grammatical relevance.
      \vskip0.6em
      A complete Transformer stacks multiple self-attention layers. Earlier
      layers might focus on basic syntax; deeper layers integrate that to grasp
      \gterm{sentiment, context, and thematic links} across the entire input.
      \vskip0.4em
      \begin{tdcode}
\begin{Verbatim}[fontsize=\scriptsize]
O(N^2 . S . D)
  N = tokens in the context
  S = number of self-attention layers
  D = number of heads per layer
\end{Verbatim}
      \end{tdcode}
    \end{tdmain}
    \begin{tdmargin}
      Transformers contain hundreds of billions or even trillions of parameters.
      This course has generally recommended fewer parameters over more, because
      fewer uses fewer resources. However, research shows that Transformers with
      more parameters \textbf{consistently outperform} those with fewer.
    \end{tdmargin}
  \end{withmargin}
\end{frame}

\begin{frame}[fragile]{How LLMs are trained}
  \begin{withmargin}
    \begin{tdmain}
      The primary ingredient is a phenomenal amount of training text, typically
      somewhat filtered. The first phase is usually unsupervised learning on
      \gterm{masked predictions}: certain tokens are intentionally hidden and
      the model trains by trying to predict them.
      \vskip0.4em
      \begin{tdcode}
\begin{Verbatim}[fontsize=\scriptsize]
Oranges are traditionally ___ by hand.
Once clipped from a tree, ___ don't ripen.
\end{Verbatim}
      \end{tdcode}
      \vskip0.4em
      \footnotesize
      Extensive training on enormous numbers of masked examples enables an LLM
      to learn that "harvested" or "picked" are high probability matches for the
      first token and "oranges" or "they" for the second.
    \end{tdmain}
    \begin{tdmargin}
      An LLM is just a neural net, so loss, the number of masked tokens the
      model correctly considered, guides the degree to which backpropagation
      updates parameter values.
      \par\medskip
      An optional further training step called \textbf{instruction tuning} can
      improve an LLM's ability to follow instructions.
      \par\medskip
      You probably will never train an LLM from scratch.
    \end{tdmargin}
  \end{withmargin}
\end{frame}

\begin{frame}[fragile]{How an LLM generates text}
  \begin{withmargin}
    \begin{tdmain}
      LLMs are essentially \gterm{autocomplete} mechanisms that can
      automatically predict thousands of tokens. You can think of a user's
      question to an LLM as the given sentence followed by an imaginary mask.
      \vskip0.4em
      \begin{tdcode}
\begin{Verbatim}[fontsize=\scriptsize]
User's question: What is the easiest trick
                 to teach a dog?
LLM's response:  ___
\end{Verbatim}
      \end{tdcode}
      \vskip0.4em
      \footnotesize
      An LLM trained on a massive number of mathematical word problems can give
      the appearance of doing sophisticated mathematical reasoning. However,
      those LLMs are basically just autocompleting a word problem prompt.
    \end{tdmain}
    \begin{tdmargin}
      For "My dog, Max, knows how to perform many traditional dog tricks", an
      LLM generates probabilities for the masked follow-on sentence:
      \par\medskip
      3.1\% "For example, he can sit, stay, and roll over."
      \par\smallskip
      2.9\% "For example, he knows how to sit, stay, and roll over."
    \end{tdmargin}
  \end{withmargin}
\end{frame}

\begin{frame}{Benefits and problems}
  \begin{withmargin}
    \begin{tdmain}
      \textbf{Benefits.} LLMs can generate clear, easy-to-understand text for a
      wide variety of target audiences, and can make predictions on tasks they
      are explicitly trained on. Some researchers claim LLMs can also make
      predictions for input they were not explicitly trained on, but other
      researchers have refuted this claim.
      \vskip0.7em
      \textbf{Problems in training.} Gathering an enormous training set.
      Consuming multiple months and enormous computational resources and
      electricity. Solving parallelism challenges.
      \vskip0.5em
      \textbf{Problems in inference.} LLMs \alertbf{hallucinate}, meaning their
      predictions often contain mistakes. They consume enormous amounts of
      computational resources and electricity. Like all ML models, they can
      exhibit all sorts of bias.
    \end{tdmain}
    \begin{tdmargin}
      Training LLMs on larger datasets typically reduces the amount of resources
      required for inference, though the larger training sets incur more
      training resources.
    \end{tdmargin}
  \end{withmargin}
  \keyterms{decoder; embedding layer; encoder; hallucination; instruction
    tuning; large language model; neural net; parameter; self attention; TPU;
    unsupervised learning}
\end{frame}

\begin{frame}{Fine-tuning}
  \begin{withmargin}
    \begin{tdmain}
      A foundation LLM is trained on enough natural language to know a
      remarkable amount about grammar, words, and idioms, and can even write
      poetry. But its generative text output isn't a solution for other kinds of
      common ML problems, such as regression or classification. For those, a
      foundation LLM serves as a \gterm{platform rather than a solution}.
      \vskip0.8em
      Fine-tuning trains on examples specific to the task your application will
      perform. Engineers can sometimes fine-tune a foundation LLM on just a few
      hundred or a few thousand training examples.
    \end{tdmain}
    \begin{tdmargin}
      Standard fine-tuning is often computationally expensive, because it
      updates the weight and bias of every parameter on each backpropagation
      iteration. \textbf{Parameter-efficient tuning} adjusts only a subset.
      \par\medskip
      A fine-tuned model contains the \textbf{same} number of parameters as the
      foundation LLM. Ten billion in, ten billion out.
    \end{tdmargin}
  \end{withmargin}
\end{frame}

\begin{frame}{Distillation}
  \begin{withmargin}
    \begin{tdmain}
      Distillation creates a smaller version of an LLM. The distilled LLM
      generates predictions much faster and requires fewer computational and
      environmental resources. However, its predictions are generally
      \gterm{not quite as good}.
      \vskip0.8em
      The most common form uses bulk inference to label data. That labeled data
      then trains a new, smaller \textbf{student} model that can be more
      affordably served. The labeled data serves as a channel by which the
      larger \textbf{teacher} model funnels its knowledge to the smaller model.
    \end{tdmain}
    \begin{tdmargin}
      For an online toxicity scorer, use a large offline scorer to label training
      data, then distill a model small enough to handle live traffic.
      \par\medskip
      A teacher can funnel a numerical score instead of a binary label. A
      numerical score is a richer training signal, enabling the student to
      predict borderline classes too.
    \end{tdmargin}
  \end{withmargin}
\end{frame}

\begin{frame}[fragile]{Prompt engineering}
  \begin{withmargin}
    \begin{tdmain}
      Humans learn well from examples. So do LLMs. Showing one example is
      \textbf{one-shot prompting}; showing several is \textbf{few-shot
      prompting}; showing none is \textbf{zero-shot prompting}.
      \vskip0.4em
      \begin{tdcode}
\begin{Verbatim}[fontsize=\scriptsize]
one-shot:   peach: drupe
            apple: ______      ->  apple: pome

few-shot:   plum: drupe
            pear: pome
            lemon: ____

zero-shot:  apple: _______
\end{Verbatim}
      \end{tdcode}
    \end{tdmain}
    \begin{tdmargin}
      Without context, the zero-shot prompt might return information about the
      \gterm{technology company} rather than the fruit.
      \par\medskip
      Prompt engineering doesn't alter the model's parameters. Prompts leverage
      the pattern-recognition abilities of the existing LLM.
    \end{tdmargin}
  \end{withmargin}
\end{frame}

\begin{frame}{Offline inference}
  \begin{withmargin}
    \begin{tdmain}
      The number of parameters in an LLM is sometimes so large that online
      inference is too slow to be practical for real-world tasks like regression
      or classification. Many engineering teams rely on \gterm{offline
      inference}, also known as bulk or static inference: rather than responding
      to queries at serving time, the trained model makes predictions in advance
      and caches them.
      \vskip0.8em
      It doesn't matter if it takes a long time for an LLM to complete its task
      if the LLM only has to perform the task once a week or once a month.
    \end{tdmain}
    \begin{tdmargin}
      Google Search used an LLM to perform offline inference in order to cache a
      list of over 800 synonyms for Covid vaccines in more than 50 languages.
      Search then used the cached list to identify queries about vaccines in
      live traffic.
      \par\medskip
      Like any form of ML, LLMs generally share the biases of the data they were
      trained on and the data they were distilled on.
    \end{tdmargin}
  \end{withmargin}
  \keyterms{backpropagation; distillation; few-shot prompting; fine-tuning;
    offline inference; one shot prompting; online inference; parameter efficient
    tuning; prompt engineering; zero-shot prompting}
\end{frame}

% ======================================================================
\sectionsubtitle{The model is about 5 percent of the codebase.}
\section[Production ML]{Production ML Systems}

\begin{frame}{\modulehead{Production ML systems}{70 minutes}}
  \begin{withmargin}
    \begin{tdmain}
      Real-world production ML systems are large ecosystems and the model is
      just a single, relatively small part. At the heart of a real-world system
      is the ML model code, but it often represents only \alertbf{5\% or less}
      of the total codebase.
      \vskip0.8em
      That's not a misprint. An ML production system devotes considerable
      resources to the input data: collecting it, verifying it, and extracting
      features from it.
    \end{tdmain}
    \begin{tdmargin}
      The other components: data collection, feature extraction, process
      management tools, data verification, configuration, machine resource
      management, monitoring, and serving infrastructure.
      \par\medskip
      The ML model code part of the course's diagram is dwarfed by the other
      nine components.
    \end{tdmargin}
  \end{withmargin}
  \slidesources{Sculley2015}
\end{frame}

\begin{frame}{Static versus dynamic training}
  \begin{withmargin}
    \begin{tdmain}
      \footnotesize
      \begin{tabular}{@{}p{0.16\linewidth}p{0.38\linewidth}p{0.38\linewidth}@{}}
        \toprule
         & \textbf{Static training (offline)} & \textbf{Dynamic training (online)}\\
        \midrule
        & Train once, serve that same model for a while
        & Train continuously or frequently, serve the most recent model\\[0.4em]
        Advantages & Simpler. You only need to develop and test the model once.
        & More adaptable. Keeps up with changes to the relationship between
          features and labels.\\[0.4em]
        Disadvantages & Sometimes staler. If the relationship changes over time,
          predictions degrade.
        & More work. You must build, test, and release a new product all the
          time.\\
        \bottomrule
      \end{tabular}
    \end{tdmain}
    \begin{tdmargin}
      If your dataset truly isn't changing over time, choose static training
      because it is cheaper. However, datasets tend to change, even those with
      features you think are as constant as sea level.
      \par\medskip
      \gterm{Even with static training, you must still monitor your input data
      for change.} A flower-buying model trained on July and August data made
      terrible predictions around Valentine's Day.
    \end{tdmargin}
  \end{withmargin}
  \keyterms{dynamic training; static training}
\end{frame}

\begin{frame}{Static versus dynamic inference}
  \begin{withmargin}
    \begin{tdmain}
      \textbf{Static inference}, also called offline or batch inference, means
      the model makes predictions on a bunch of common unlabeled examples and
      caches them somewhere. \textbf{Dynamic inference}, also called online or
      real-time inference, means the model only makes predictions on demand.
      \vskip0.7em
      \footnotesize
      A very complex model that takes one hour to infer a prediction would be an
      excellent situation for static inference. If it mistakenly used dynamic
      inference and many clients requested predictions around the same time,
      most of them wouldn't receive a prediction for \gterm{hours or days}. A
      model that infers in 2 milliseconds can serve clients quickly and
      efficiently through dynamic inference.
    \end{tdmain}
    \begin{tdmargin}
      {\color{tdfg}\textbf{Static inference.}} Don't need to worry much about
      cost of inference; can do post-verification before pushing. But can only
      serve cached predictions, and update latency is measured in hours or days.
      \par\medskip
      {\color{tdfg}\textbf{Dynamic inference.}} Can infer on any new item as it
      comes in, great for long tail predictions. But compute intensive and
      latency sensitive, which may limit model complexity, and monitoring needs
      are more intensive.
    \end{tdmargin}
  \end{withmargin}
  \keyterms{example; inference; unlabeled example}
\end{frame}

\begin{frame}{When to transform data}
  \begin{withmargin}
    \begin{tdmain}
      \textbf{Before training.} Write code or use specialized tools to transform
      the raw data, then store it where the model can ingest it. The system
      transforms raw data only once, and can analyze the entire dataset to
      determine the best transformation strategy. But you must recreate the
      transformations at prediction time. \alertbf{Beware of training-serving
      skew.}
      \vskip0.6em
      \textbf{While training.} The transformation is part of the model code. You
      can still use the same raw data files if you change the transformations,
      and you're ensured the same transformations at training and prediction
      time. But complicated transforms increase model latency, and
      transformations occur for each and every batch.
    \end{tdmain}
    \begin{tdmargin}
      Training-serving skew is more dangerous with dynamic inference, because
      the software that transforms the raw dataset usually differs from the
      software that serves predictions. Systems using static inference can
      sometimes use the same software.
      \par\medskip
      Per-batch transformation is tricky: a Z-score of $-2.5$ in one batch won't
      mean the same as $-2.5$ in another. Precompute the mean and standard
      deviation across the entire dataset and use them as constants.
    \end{tdmargin}
  \end{withmargin}
  \keyterms{training-serving skew; Z-score normalization}
\end{frame}

\begin{frame}{Testing a deployment}
  \begin{withmargin}
    \begin{tdmain}
      If only deploying a model were as easy as pressing a big Deploy button. A
      full machine learning system requires tests for validating input data,
      validating feature engineering, validating the quality of new model
      versions, validating serving infrastructure, and testing integration
      between pipeline components.
      \vskip0.8em
      Many software engineers favor test-driven development. However, TDD can be
      tricky in ML. Before training your model, you \gterm{can't write a test to
      validate the loss}. Instead, you must first discover the achievable loss
      during model development and then test new model versions against it.
    \end{tdmain}
    \begin{tdmargin}
      {\color{tdfg}\textbf{Reproducible training.}}
      \par\medskip
      Deterministically seed the random number generator. Initialize model
      components in a fixed order. Take the average of several runs. Use version
      control, even for preliminary iterations.
      \par\medskip
      Even after following these guidelines, other sources of nondeterminism
      might still exist.
    \end{tdmargin}
  \end{withmargin}
\end{frame}

\begin{frame}{Two kinds of quality degradation}
  \begin{withmargin}
    \begin{tdmain}
      \textbf{Sudden degradation.} A bug in the new version could cause
      significantly lower quality. Validate new versions by checking their
      quality against the previous version.
      \vskip0.7em
      \textbf{Slow degradation.} Your test for sudden degradation might not
      detect a slow decline over multiple versions. Ensure your model's
      predictions on a validation dataset meet a \gterm{fixed threshold}. If
      your validation dataset deviates from live data, update it and ensure the
      model still meets the same threshold.
      \vskip0.6em
      \footnotesize
      If your model is updated faster than your server, it could have different
      software dependencies from the server. Stage the model in a sandboxed
      version of the server.
    \end{tdmain}
    \begin{tdmargin}
      To test updates to API calls without retraining, write a unit test that
      generates random input data and runs a single step of gradient descent. If
      it completes without errors, then any updates to the API probably haven't
      ruined your model.
      \par\medskip
      To run integration tests faster, train on a subset of the data or with a
      simpler model.
    \end{tdmargin}
  \end{withmargin}
\end{frame}

\begin{frame}{Write a data schema}
  \begin{withmargin}
    \begin{tdmain}
      To monitor your data, continuously check it against expected statistical
      values by writing rules that the data must satisfy. This collection of
      rules is a \gterm{data schema}.
      \vskip0.6em
      \begin{itemize}
        \item Understand the range and distribution of your features. For
          categorical features, understand the set of possible values.
        \item Encode your understanding into the schema. Ensure user-submitted
          ratings are always in the range 1 to 5. Check that the word "the"
          occurs most frequently in an English text feature.
        \item Test your data against the schema, to catch anomalies, unexpected
          values of categorical variables, and unexpected data distributions.
      \end{itemize}
    \end{tdmain}
    \begin{tdmargin}
      Your model doesn't train on raw data; it trains on feature engineered
      data. Check that separately with unit tests: all numeric features scaled
      between 0 and 1, one-hot vectors containing a single 1 and $N-1$ zeroes,
      Z-score means of 0, outliers handled.
    \end{tdmargin}
  \end{withmargin}
\end{frame}

\begin{frame}{Check metrics for important data slices}
  \begin{withmargin}
    \begin{tdmain}
      A successful whole sometimes obscures an unsuccessful subset. A model with
      great overall metrics might still make terrible predictions for certain
      situations.
      \vskip0.7em
      \begin{tdblock}
        Your unicorn model performs well overall, but performs poorly when
        making predictions for the Sahara desert.
      \end{tdblock}
      \vskip0.6em
      \footnotesize
      If you are the kind of engineer satisfied with an overall great AUC, you
      might not notice. Subsets of data like this are called \gterm{data
      slices}. Compare model metrics for these slices against the metrics for
      your entire dataset.
    \end{tdmain}
    \begin{tdmargin}
      Model metrics don't necessarily measure the real-world impact of your
      model. Changing a hyperparameter might increase AUC, but how did the
      change affect user experience?
      \par\medskip
      To measure real-world impact you need separate metrics. You could survey
      users of your model to confirm that they really did see a unicorn when the
      model predicted they would.
    \end{tdmargin}
  \end{withmargin}
\end{frame}

\begin{frame}{Training-serving skew}
  \begin{withmargin}
    \begin{tdmain}
      \footnotesize
      \begin{tabular}{@{}p{0.14\linewidth}p{0.26\linewidth}p{0.24\linewidth}p{0.26\linewidth}@{}}
        \toprule
        \textbf{Type} & \textbf{Definition} & \textbf{Example} & \textbf{Solution}\\
        \midrule
        Schema skew & Training and serving input data do not conform to the same
          schema & The format or distribution of serving data changes while the
          model continues to train on old data & Use the same schema to validate
          both; separately check statistics the schema misses, such as the
          fraction of missing values\\[0.4em]
        Feature skew & Engineered data differs between training and serving
          & Feature engineering code differs between training and serving
          & Apply the same statistical rules across both; track the number of
          skewed features and the ratio of skewed examples per feature\\
        \bottomrule
      \end{tabular}
    \end{tdmain}
    \begin{tdmargin}
      {\color{tdfg}\textbf{The Golden Rule.}}
      \par\medskip
      Ensure that training and production \gterm{mimic each other} as closely as
      possible. The more closely your training task matches your prediction
      task, the better your ML system will perform.
      \par\medskip
      Always consider what data is available to your model at prediction time.
    \end{tdmargin}
  \end{withmargin}
\end{frame}

\begin{frame}{Label leakage}
  \begin{withmargin}
    \begin{tdmain}
      Label leakage means the ground truth labels you're trying to predict have
      inadvertently entered your training features. Label leakage is sometimes
      very difficult to detect.
      \vskip0.8em
      A cancer model using patient age, gender, prior conditions,
      \gterm{hospital name}, vital signs, test results, and heredity performs
      exceedingly well on validation and test sets and terribly in the real
      world. Certain hospitals specialize in treating cancer, so the model
      learned that patients assigned to those hospitals very likely had cancer.
      At inference time most patients hadn't yet been assigned to a hospital.
    \end{tdmain}
    \begin{tdmargin}
      Another example: predicting daily revenue using the number of customers as
      a feature. You don't know the number of customers at prediction time,
      before the day's sales are complete.
      \par\medskip
      When you get amazing evaluation metrics, like 0.99 AUC, look for these
      sorts of features that can bleed into your label.
    \end{tdmargin}
  \end{withmargin}
  \slidesources{Kaufman2011}
\end{frame}

\begin{frame}{Monitor model age and performance}
  \begin{withmargin}
    \begin{tdmain}
      If the serving data evolves with time but your model isn't retrained
      regularly, you will see a decline in model quality. Track the time since
      the model was retrained on new data and set a threshold age for alerts.
      Monitor the model's age \gterm{throughout the pipeline} to catch pipeline
      stalls.
      \vskip0.6em
      \footnotesize
      \begin{itemize}
        \item Track model performance by versions of code, model, and data.
        \item Test training steps per second against the previous version and a
          fixed threshold.
        \item Catch memory leaks by setting a threshold for memory use.
        \item Monitor API response times and track their percentiles.
        \item Monitor the number of queries answered per second.
      \end{itemize}
    \end{tdmain}
    \begin{tdmargin}
      During training, your weights and layer outputs should not be
      \texttt{NaN} or \texttt{Inf}. Write tests to check. Additionally, test
      that more than half of the outputs of a layer are not zero.
    \end{tdmargin}
  \end{withmargin}
\end{frame}

\begin{frame}{Testing quality on served data}
  \begin{withmargin}
    \begin{tdmain}
      Testing quality in serving is hard because real-world data is not always
      labeled. If your serving data is not labeled, consider these tests.
      \vskip0.6em
      \begin{itemize}
        \item Generate labels using human raters.
        \item Investigate models that show significant statistical bias in
          predictions.
        \item Track real-world metrics. If you're classifying spam, compare your
          predictions to user-reported spam.
        \item Serve a new model version on a \gterm{fraction} of your queries.
          As you validate it, gradually switch all queries to the new version.
      \end{itemize}
      \vskip0.4em
      \footnotesize
      Remember to monitor both sudden and slow degradation in prediction
      quality.
    \end{tdmain}
    \begin{tdmargin}
      In production, calculating a model's quality is different from calculating
      it during training. You won't necessarily have access to the ground truth.
      Write custom monitoring instrumentation to capture metrics that act as a
      proxy: in a mail app, monitor the percentage of mail users move to spam.
      If it jumps from 0.5\% to 3\%, that signals a potential issue.
    \end{tdmargin}
  \end{withmargin}
\end{frame}

\begin{frame}{Randomization and hashing}
  \begin{withmargin}
    \begin{tdmain}
      Make your data generation pipeline \gterm{reproducible}. For a fair
      experiment, your datasets should be identical except for the new feature
      you are testing.
      \vskip0.6em
      \textbf{Seed your random number generators}, so the RNG outputs the same
      values in the same order each time, recreating your dataset.
      \vskip0.5em
      \textbf{Use invariant hash keys.} The inputs to your hash function
      shouldn't change each time you run the data generation program. Don't use
      the current time or a random number in your hash if you want to recreate
      your hashes on demand.
    \end{tdmain}
    \begin{tdmargin}
      If the hash key only used the query, then across multiple days you'd
      always include that query or always exclude it. That is bad: the training
      set sees a less diverse set of queries, and evaluation sets become
      artificially hard because they won't overlap with training data.
      \par\medskip
      Instead hash on \textbf{query plus query date}, giving a different hashing
      each day. Use the date the query was logged, not the current date, so
      hashing stays deterministic.
    \end{tdmargin}
  \end{withmargin}
  \keyterms{AUC; clipping; ground truth; hashing; label leakage; one-hot
    encoding; outliers; scaling; training-serving skew}
\end{frame}

\begin{frame}{Questions to ask in production}
  \begin{withmargin}
    \begin{tdmain}
      \textbf{Is each feature helpful?} Continuously monitor your model to
      remove features that contribute little or nothing. If the input data for
      that feature abruptly changes, your model's behavior might also abruptly
      change in undesirable ways. Also ask whether the usefulness of the feature
      justifies the \gterm{maintenance burden} of including it.
      \vskip0.7em
      \textbf{Is your data source reliable?} Is the signal coming from a server
      that crashes under heavy load? From humans that go on vacation every
      August? Does the system that computes your model's input data ever change?
      How often, and how will you know?
    \end{tdmain}
    \begin{tdmargin}
      It is always tempting to add more features. A new feature that makes
      predictions slightly better certainly seems better than slightly worse
      predictions; however, the extra feature adds to your maintenance burden.
      \par\medskip
      Consider creating your own copy of the data you receive from the upstream
      process. Then only advance to the next version of the upstream data when
      you are certain that it is safe to do so.
    \end{tdmargin}
  \end{withmargin}
\end{frame}

\begin{frame}{Is your model part of a feedback loop?}
  \begin{withmargin}
    \begin{tdmain}
      Sometimes a model can affect its own training data: the results from some
      models become, directly or indirectly, input features to that same model.
      Sometimes a model can affect another model.
      \vskip0.7em
      \footnotesize
      Model A is a bad predictive model for stock prices. Since Model A is
      buggy, it mistakenly decides to buy stock in Stock X. Those purchases
      \gterm{drive up the price} of Stock X. Model B uses the price of Stock X
      as an input feature, so Model B can come to false conclusions about its
      value. Model B's behavior, in turn, can affect Model A, possibly
      triggering a tulip mania or a slide in Company X's stock.
    \end{tdmain}
    \begin{tdmargin}
      {\color{tdfg}\textbf{Susceptible.}} A traffic-forecasting model using
      beach crowd size. A book-recommendation model using popularity. A
      university-ranking model that rates schools by selectivity.
      \par\medskip
      {\color{tdfg}\textbf{Not susceptible.}} An election-results model that
      forecasts after the polls have closed. A housing-value model using size,
      bedrooms, and location.
    \end{tdmargin}
  \end{withmargin}
\end{frame}

% ======================================================================
\sectionsubtitle{Automating the repetitive parts of the workflow.}
\section[AutoML]{AutoML}

\begin{frame}{\modulehead{AutoML}{30 minutes}}
  \begin{withmargin}
    \begin{tdmain}
      With manual training, you write code using an ML framework to create a
      model, choose which algorithms to explore, and iteratively tune
      hyperparameters. Two things make this heavy.
      \vskip0.6em
      \textbf{Repetitive tasks.} During model development you typically need to
      explore different combinations of algorithms and hyperparameters. For
      small or exploratory projects this may not be a problem, but for larger
      projects these repetitive tasks can be time consuming.
      \vskip0.5em
      \textbf{Specialized skills.} If a team does not have a dedicated data
      scientist, doing this work manually might not even be \gterm{feasible}.
    \end{tdmain}
    \begin{tdmargin}
      AutoML is a process of automating certain tasks in a machine learning
      workflow: a set of tools and technologies that make building models faster
      and more accessible to a wider group of users.
      \par\medskip
      The tasks: feature engineering and feature selection; identifying an
      appropriate ML algorithm and selecting the best hyperparameters;
      evaluating metrics generated during training.
    \end{tdmargin}
  \end{withmargin}
\end{frame}

\begin{frame}{AutoML tools and workflow}
  \begin{withmargin}
    \begin{tdmain}
      \textbf{No-code tools} typically take the form of web applications that
      let you configure and run experiments through a user interface.
      \textbf{API and CLI tools} provide advanced automation features, but
      require more, sometimes significantly more, programming and ML expertise.
      \vskip0.7em
      \footnotesize
      The high level workflow steps are the same as for custom training; the
      main difference is that AutoML handles some tasks for you. Before
      importing data you still need to \gterm{label your data}, clean and format
      it, and perform feature transformations the tool does not support.
    \end{tdmain}
    \begin{tdmargin}
      Real-world data tends to be messy, so expect to clean your data before
      using it. Even with AutoML you need to determine the best treatments for
      your particular dataset and problem.
      \par\medskip
      Ensure that the tool you choose supports the objectives of your project,
      your data source, the data types in your dataset, and the size of your
      dataset.
    \end{tdmargin}
  \end{withmargin}
\end{frame}

\begin{frame}{Configuring an AutoML run}
  \begin{withmargin}
    \begin{tdmain}
      \textbf{Import} your data by specifying the data source. During import,
      the tool assigns a semantic data type to each data value.
      \vskip0.5em
      \textbf{Analyze} and \textbf{refine}. During import, AutoML tools try to
      determine the correct semantic type for each feature, but these are only
      \gterm{guesses}. Check the types designated to all features and change
      them if assigned incorrectly.
      \vskip0.5em
      \footnotesize
      \textbf{Configure.} Select the ML problem type. Select which column is the
      label. Select the set of features. Select the set of ML algorithms AutoML
      considers in the model search. Select the evaluation metric AutoML uses to
      choose the best model.
    \end{tdmain}
    \begin{tdmargin}
      Postal codes stored as numbers would be detected as continuous numeric
      data. That is incorrect, and you would want to change the semantic type to
      categorical.
      \par\medskip
      Some tools allow custom transformations. For a housing dataset with a
      \texttt{description} column, a \texttt{LENGTH} function could generate a
      new \texttt{description\_length} feature.
      \par\medskip
      Training may take a while to complete, on the order of hours.
    \end{tdmargin}
  \end{withmargin}
\end{frame}

\begin{frame}{Benefits of AutoML}
  \begin{withmargin}
    \begin{tdmain}
      \footnotesize
      \begin{itemize}
        \item \textbf{Save time} by avoiding extensive manual experimentation to
          find the best model.
        \item \textbf{Improve quality.} AutoML tools can comprehensively search
          for the highest quality model.
        \item \textbf{Build without specialized skills.} A side effect of
          automating ML tasks is that it \gterm{democratizes} ML.
        \item \textbf{Smoke test a dataset.} Even for an expert, AutoML quickly
          gives a baseline estimate for whether a dataset has enough signal in
          all of its noise. If AutoML can't build even a mediocre model, it
          might not be worth building a good one by hand.
        \item \textbf{Evaluate a dataset.} Even if you don't use the resulting
          model, AutoML may help you determine which features aren't worth
          gathering.
        \item \textbf{Enforce best practices} on each model search.
      \end{itemize}
    \end{tdmain}
    \begin{tdmargin}
      \vfill
      After training, examine feature importance metrics, the architecture and
      hyperparameters used to build the model, and the plots and metrics
      collected during training.
      \vfill
    \end{tdmargin}
  \end{withmargin}
\end{frame}

\begin{frame}{Limitations of AutoML}
  \begin{withmargin}
    \begin{tdmain}
      \footnotesize
      \begin{itemize}
        \item \textbf{Model quality may not be as good as manual training.} A
          motivated expert with enough time can often create a model with better
          prediction quality.
        \item \textbf{Model search and complexity can be opaque.} It is
          difficult to have insight into how the tool arrived at the best model,
          and models generated with AutoML are difficult to reproduce manually.
        \item \textbf{Multiple runs may show more variance.} Different runs may
          search different portions of the space and wind up in moderately, or
          even significantly, different places.
        \item \textbf{Models can't be customized during training.} If your use
          case requires customization or tweaking during the training process,
          AutoML may not be the right choice.
      \end{itemize}
    \end{tdmain}
    \begin{tdmargin}
      One thing you can count on when building a model from scratch is that you
      need \gterm{large amounts of data}. The advantage with AutoML is that you
      can mostly ignore the architecture and hyperparameter search and focus
      primarily on the quality of your data.
      \par\medskip
      Specialized AutoML systems using transfer learning can train with
      significantly less data: a few hundred labeled images rather than hundreds
      of thousands.
    \end{tdmargin}
  \end{withmargin}
\end{frame}

\begin{frame}{Is AutoML right for your project?}
  \vfill
  \begin{takeaway}
    AutoML for a team with limited ML experience,\\
    or an experienced team seeking productivity gains\\
    with no customization requirements.
  \end{takeaway}
  \vfill
  {\small Custom manual training is more appropriate when model quality is
  important and the team needs to customize their model. It may require more
  time for experimentation, but the team can often achieve a
  \alertbf{higher quality} model.}
  \keyterms{AutoML; transfer learning}
\end{frame}

% ======================================================================
\sectionsubtitle{Auditing data and predictions for bias.}
\section[Fairness]{ML Fairness}

\begin{frame}{\modulehead{ML fairness}{110 minutes}}
  \begin{withmargin}
    \begin{tdmain}
      Evaluating a machine learning model responsibly requires doing more than
      just calculating overall loss metrics. Before putting a model into
      production, it's critical to \alertbf{audit training data and evaluate
      predictions for bias}.
      \vskip0.8em
      ML models are not inherently objective. Practitioners train models by
      feeding them a dataset of training examples, and human involvement in the
      provision and curation of this data can make a model's predictions
      susceptible to bias.
    \end{tdmain}
    \begin{tdmargin}
      {\color{tdfg}\textbf{Objectives.}}
      \par\medskip
      Become aware of common human biases that can inadvertently be reproduced
      by ML algorithms. Proactively explore data to identify sources of bias
      before training a model. Evaluate model predictions for bias.
      \par\medskip
      Wikipedia's catalog of cognitive biases enumerates over 100 different
      types of human bias that can affect our judgment.
    \end{tdmargin}
  \end{withmargin}
\end{frame}

\begin{frame}{Reporting, historical, and automation bias}
  \begin{withmargin}
    \begin{tdmain}
      \footnotesize
      \textbf{Reporting bias} occurs when the frequency of events, properties,
      or outcomes captured in a dataset does not accurately reflect their
      real-world frequency. People tend to focus on documenting circumstances
      that are unusual or especially memorable.
      \vskip0.5em
      \textbf{Historical bias} occurs when historical data reflects inequities
      that existed in the world at that time.
      \vskip0.5em
      \textbf{Automation bias} is a tendency to favor results generated by
      automated systems over those generated by non-automated systems,
      \gterm{irrespective of the error rates of each}.
    \end{tdmain}
    \begin{tdmargin}
      A sentiment model trained on book reviews sees mostly extreme opinions,
      because people were less likely to review a book they did not respond to
      strongly. It is then less able to predict sentiment of reviews using more
      subtle language.
      \par\medskip
      A 1960s city housing dataset contains home-price data reflecting
      discriminatory lending practices in effect that decade.
      \par\medskip
      A sprocket manufacturer was eager to deploy a groundbreaking tooth-defect
      model, until the factory supervisor pointed out its precision and recall
      were both 15\% lower than human inspectors.
    \end{tdmargin}
  \end{withmargin}
\end{frame}

\begin{frame}{Selection bias}
  \begin{withmargin}
    \begin{tdmain}
      Selection bias occurs if a dataset's examples are chosen in a way that is
      not reflective of their real-world distribution.
      \vskip0.6em
      \footnotesize
      \textbf{Coverage bias} occurs if data is not selected in a representative
      fashion. Consumers who bought a competing product were never surveyed, so
      that group was not represented in the training data.
      \vskip0.4em
      \textbf{Non-response bias}, also participation bias, occurs if data ends
      up unrepresentative due to participation gaps. Consumers who bought the
      competing product were 80\% more likely to refuse to complete the survey.
      \vskip0.4em
      \textbf{Sampling bias} occurs if proper randomization is not used. The
      surveyor chose the first 200 consumers who responded to an email, who
      might have been \gterm{more enthusiastic} than average purchasers.
    \end{tdmain}
    \begin{tdmargin}
      All three examples come from the same scenario: a model trained to predict
      future sales of a new product based on phone surveys.
    \end{tdmargin}
  \end{withmargin}
\end{frame}

\begin{frame}{Group attribution bias}
  \begin{withmargin}
    \begin{tdmain}
      Group attribution bias is a tendency to generalize what is true of
      individuals to the entire group to which they belong.
      \vskip0.7em
      \textbf{In-group bias} is a preference for members of your own group, or
      for characteristics that you also share. Two practitioners training a
      resume-screening model are predisposed to believe that applicants who
      attended the same computer-science academy as they both did are
      \gterm{more qualified}.
      \vskip0.6em
      \textbf{Out-group homogeneity bias} is a tendency to stereotype individual
      members of a group to which you do not belong, or to see their
      characteristics as more uniform. The same two practitioners believe that
      all applicants who did not attend a computer-science academy lack
      sufficient expertise.
    \end{tdmain}
    \begin{tdmargin}
      \textbf{Implicit bias} occurs when assumptions are made based on one's own
      model of thinking and personal experiences that don't necessarily apply
      more generally.
      \par\medskip
      A practitioner training a gesture-recognition model uses a head shake as a
      feature to indicate "no". However, in some regions of the world, a head
      shake actually signifies "yes".
    \end{tdmargin}
  \end{withmargin}
\end{frame}

\begin{frame}{Confirmation and experimenter's bias}
  \begin{withmargin}
    \begin{tdmain}
      \textbf{Confirmation bias} occurs when model builders unconsciously
      process data in ways that affirm pre-existing beliefs and hypotheses.
      \vskip0.7em
      \textbf{Experimenter's bias} occurs when a model builder keeps training a
      model until it produces a result that \gterm{aligns with their original
      hypothesis}.
      \vskip0.7em
      \footnotesize
      Both are illustrated with the same practitioner: someone building a model
      that predicts aggressiveness in dogs, who had an unpleasant encounter with
      a hyperactive toy poodle as a child and ever since has associated the
      breed with aggression.
    \end{tdmain}
    \begin{tdmargin}
      Under confirmation bias, the practitioner unconsciously discarded features
      that provided evidence of docility in smaller dogs.
      \par\medskip
      Under experimenter's bias, when the trained model predicted most toy
      poodles to be relatively docile, the practitioner retrained the model
      several more times until it produced a result showing smaller poodles to
      be more violent.
    \end{tdmargin}
  \end{withmargin}
  \keyterms{automation bias; confirmation bias; coverage bias; experimenter's
    bias; group attribution bias; historical bias; implicit bias; in-group bias;
    out-group homogeneity bias; non-response bias; reporting bias; sampling
    bias; selection bias}
\end{frame}

\begin{frame}{Red flags in your dataset}
  \begin{withmargin}
    \begin{tdmain}
      \textbf{Missing feature values.} If one or more features have missing
      values for a large number of examples, that could indicate that certain
      key characteristics of your dataset are under-represented.
      \vskip0.7em
      \textbf{Unexpected feature values.} Look for examples that contain feature
      values that stand out as especially uncharacteristic or unusual. These
      could indicate problems during data collection.
      \vskip0.6em
      \textbf{Data skew.} Any sort of skew, where certain groups may be under-
      or over-represented relative to their real-world prevalence, can introduce
      bias. Break out results \gterm{by subgroup}, not just in aggregate.
    \end{tdmain}
    \begin{tdmargin}
      In a rescue-dog adoptability dataset, 1{,}500 of 5{,}000 examples are
      missing temperament values. If that correlates with breed or with age,
      predictions become less accurate for those groups.
      \par\medskip
      Even missing for all dogs from big cities is a risk: location may serve as
      a proxy for physical characteristics.
      \par\medskip
      A labrador retriever recorded as 35 years old is implausible: the oldest
      verified dog, Bluey, lived 29 years and 5 months.
    \end{tdmargin}
  \end{withmargin}
\end{frame}

\begin{frame}{Mitigating bias}
  \begin{withmargin}
    \begin{tdmain}
      \textbf{Augment the training data.} If an audit uncovered missing,
      incorrect, or skewed data, the most straightforward fix is often to
      collect more. The downside is that this can be infeasible, either due to a
      lack of available data or resource constraints.
      \vskip0.7em
      \textbf{Adjust the optimization function.} Log loss does not take subgroup
      membership into consideration. Instead choose a function designed to
      penalize errors in a \gterm{fairness-aware} fashion that counteracts the
      imbalances you've identified.
    \end{tdmain}
    \begin{tdmargin}
      \textbf{MinDiff} balances the errors for two different slices of data by
      adding a penalty for differences in the prediction distributions for the
      two groups.
      \par\medskip
      \textbf{Counterfactual Logit Pairing} ensures that changing a sensitive
      attribute of a given example doesn't alter the model's prediction for that
      example.
      \par\medskip
      The two are complementary: collect more data to reduce a discrepancy by
      30\%, then use MinDiff to reduce it by another 50\%.
    \end{tdmargin}
  \end{withmargin}
  \keyterms{bias (ethics/fairness); log loss}
\end{frame}

\begin{frame}{Evaluating for bias}
  \begin{withmargin}
    \begin{tdmain}
      Metrics calculated against an entire test or validation set don't always
      give an accurate picture of how fair the model is. Great overall
      performance for a majority of examples may \alertbf{mask poor performance
      on a minority subset}. Aggregate metrics such as precision, recall, and
      accuracy are not necessarily going to expose these issues.
      \vskip0.8em
      The running example: an admissions model selects 20 students from a pool
      of 100 candidates, belonging to two demographic groups, a majority group
      of 80 students and a minority group of 20 students.
    \end{tdmain}
    \begin{tdmargin}
      There are a variety of metrics, each of which provides a different
      mathematical definition of fairness. The module explores three in depth:
      \par\medskip
      demographic parity,
      \par\smallskip
      equality of opportunity,
      \par\smallskip
      counterfactual fairness.
    \end{tdmargin}
  \end{withmargin}
\end{frame}

\begin{frame}{Demographic parity}
  \begin{withmargin}
    \begin{tdmain}
      If the two admissions rates are equal, the model's predictions exhibit
      demographic parity: a student's chance of being admitted doesn't vary by
      demographic group.
      \vskip0.5em
      \footnotesize
      \begin{tabular}{@{}lrr@{}}
        \toprule
         & \textbf{Majority group} & \textbf{Minority group}\\
        \midrule
        Accepted & 16 & 4\\
        Rejected & 64 & 16\\
        \midrule
        Acceptance rate & \gterm{20\%} & \gterm{20\%}\\
        \bottomrule
      \end{tabular}
      \vskip0.5em
      The key benefit is that both groups are represented in the admitted class
      in the same proportions as they are in the candidate pool.
    \end{tdmain}
    \begin{tdmargin}
      The significant drawback: it does not take the \textbf{distribution} of
      predictions for each group, the number classified as qualified versus
      unqualified, into account when deciding how the 20 slots should be
      allocated.
    \end{tdmargin}
  \end{withmargin}
\end{frame}

\begin{frame}{Where demographic parity breaks}
  \begin{withmargin}
    \begin{tdmain}
      Break the same pool out by qualification as well as group. Of 35 qualified
      majority students, 16 were accepted. Of 15 qualified minority students, 4
      were accepted.
      \begin{align*}
        \text{Majority acceptance rate} &= \frac{16}{35} = \gterm{46\%}\\[0.3em]
        \text{Minority acceptance rate} &= \frac{4}{15} = \gterm{27\%}
      \end{align*}
      \vskip0.2em
      \footnotesize
      Even though both groups have an overall acceptance rate of 20\%,
      satisfying demographic parity, the rates for qualified students differ
      sharply.
    \end{tdmain}
    \begin{tdmargin}
      Where the distribution of a preferred label such as "qualified" varies
      greatly between groups, demographic parity may not be the optimal metric
      for evaluating fairness.
      \par\medskip
      There may be additional benefits and drawbacks not discussed here, worth
      considering depending on the goals the model is designed to achieve and
      the social context in which its predictions will be used.
    \end{tdmargin}
  \end{withmargin}
  \keyterms{bias (ethics/fairness); demographic parity}
\end{frame}

\begin{frame}{Equality of opportunity}
  \begin{withmargin}
    \begin{tdmain}
      Compare the acceptance rates for just the \gterm{qualified} candidates. If
      those are equal, the model exhibits equality of opportunity: students with
      the preferred label have an equal chance of being admitted, irrespective
      of group.
      \vskip0.5em
      \footnotesize
      \begin{tabular}{@{}lrrrr@{}}
        \toprule
         & \multicolumn{2}{c}{\textbf{Majority group}} & \multicolumn{2}{c}{\textbf{Minority group}}\\
         & Accepted & Rejected & Accepted & Rejected\\
        \midrule
        Qualified   & 14 & 21 & 6 & 9\\
        Unqualified & 0  & 45 & 0 & 5\\
        \bottomrule
      \end{tabular}
      \vskip0.4em
      Both qualified groups have an acceptance rate of 40\%.
    \end{tdmain}
    \begin{tdmargin}
      These predictions do \textbf{not} satisfy demographic parity: a majority
      student has a 17.5\% chance of acceptance, a minority student 30\%. But a
      qualified student has a 40\% chance either way, arguably a fairer outcome
      in this use case.
      \par\medskip
      The benefit: the ratio of positive to negative predictions may vary across
      groups, provided the model is equally successful at predicting the
      preferred label.
    \end{tdmargin}
  \end{withmargin}
  \slidesources{Hardt2016}
\end{frame}

\begin{frame}{Drawbacks of equality of opportunity}
  \begin{withmargin}
    \begin{tdmain}
      It is designed for use cases where there is a \gterm{clear-cut preferred
      label}. If it's equally important that the model predict both the positive
      and the negative class for all groups, it may make sense to use
      \textbf{equalized odds}, which enforces equal success rates for both
      labels.
      \vskip0.8em
      It also assesses fairness by comparing error rates in aggregate for
      demographic groups, which may not always be feasible. If the admissions
      dataset did not have a \texttt{demographic\_group} feature, it would not
      be possible to break out acceptance rates and compare them.
    \end{tdmain}
    \begin{tdmargin}
      It is possible for predictions to satisfy \textbf{both} demographic parity
      and equality of opportunity. In the course's worked confusion matrices,
      both groups have a positive prediction rate of 25\% and a true positive
      rate of 33\%.
    \end{tdmargin}
  \end{withmargin}
\end{frame}

\begin{frame}{Counterfactual fairness}
  \begin{withmargin}
    \begin{tdmain}
      Often demographic-group membership is recorded for only a small percentage
      of examples, for instance students who opted to self-identify. With
      demographic data missing for 94\% of examples, neither demographic parity
      nor equality of opportunity is feasible.
      \vskip0.8em
      Counterfactual fairness stipulates that two examples that are
      \gterm{identical in all respects}, except a given sensitive attribute,
      should result in the same model prediction. For the 6\% of examples that
      do contain demographic features, compare pairs of individual predictions
      and see if they have been treated equitably.
    \end{tdmain}
    \begin{tdmargin}
      Because parity and equality of opportunity assess groups in aggregate,
      they may mask bias issues at the level of individual predictions.
      \par\medskip
      A model might accept qualified candidates from both groups in the same
      proportion, but reject the most qualified minority candidate while
      accepting the most qualified majority candidate with the exact same
      credentials.
    \end{tdmargin}
  \end{withmargin}
\end{frame}

\begin{frame}{Choosing a fairness metric}
  \begin{withmargin}
    \begin{tdmain}
      The key downside of counterfactual fairness is that it doesn't provide as
      holistic a view of bias. Identifying and remediating a handful of
      inequities in pairs of examples may not be sufficient to address
      \gterm{systemic} bias issues that affect entire subgroups.
      \vskip0.8em
      Some definitions of fairness are even \alertbf{mutually incompatible},
      meaning it may be impossible to satisfy them simultaneously for a given
      model's predictions. So how do you choose? Consider the context in which
      the model is being used and the overarching goals. Is the goal equal
      representation, or equal opportunity?
    \end{tdmain}
    \begin{tdmargin}
      Where possible, practitioners can do both an aggregate fairness analysis,
      using demographic parity or equality of opportunity, and a counterfactual
      fairness analysis, to gain the widest range of insights.
    \end{tdmargin}
  \end{withmargin}
  \slidesources{Barocas2019}
  \keyterms{bias (ethics/fairness); demographic parity; equality of opportunity;
    equalized odds; mutually incompatible}
\end{frame}
